\chapter{機械学習の枠組み：データから関数を学ぶ}\label{ch:ml}

\bookmap{2}

\begin{chapterlead}
\textbf{ここまで}: 第 I 部で、数学の言葉（微分・確率・交差エントロピー）と、
計算の道具（NumPy と PyTorch のテンソル）を準備しました。\\[1mm]
\textbf{この章の問い}: プログラムを書く代わりに「データから関数を作る」とは、どういうことでしょうか。
何を最小化すればよく、なぜそれでうまくいくのでしょうか。\\[1mm]
\textbf{この章を終えると}:
\begin{itemize}
\item 機械学習を「データ・モデル・損失関数・最適化」の 4 要素に分けて説明できます。
\item 線形回帰を、勾配の導出から勾配降下法の実装まで、自分の手で書けます。
\item 二乗誤差も交差エントロピーも、「最尤推定」という一つの考え方から出てくることを説明できます。
\item softmax 回帰で分類問題を解き、その勾配が「予測 $-$ 正解」になることを導けます。
\item ミニバッチ、エポック、過学習、正則化、ハイパーパラメータという言葉を、実験と結びつけて使えます。
\end{itemize}
\end{chapterlead}

\begin{readingguide}
\item[\rgpre] \chapref{ch:math}の微分と勾配（\secref{sec:math-calculus}）、交差エントロピー（\secref{sec:math-info}）、softmax（\secref{sec:math-softmax}）と、\chapref{ch:pytorch}の自動微分（\secref{sec:pt-autograd}）を思い出せれば読めます。
\item[\rgmust] \secref{sec:ml-rules-vs-examples}〜\secref{sec:ml-autograd}：機械学習の 4 要素と、勾配降下法による線形回帰の学習。ここで作る学習のループを、本書の最後まで使います。
\item[\rgskim] \secref{sec:ml-mle}：二乗誤差を選ぶ理由。節末の囲みの「損失は、負の対数尤度として作る」というレシピだけ押さえれば、先へ進めます。
\item[\rgmust] \secref{sec:ml-classification}〜\secref{sec:ml-generalization}：交差エントロピー損失、ミニバッチ、過学習。言語モデルの損失と学習の進め方は、ここの言葉でできています。
このうち、勾配の導出、実装の細部、過学習の対策の三つの小節は、流し読みで構いません（見出しの下に、流し読みの印があります）。
\item[\rgskim] \secref{sec:ml-map}：学習の種類の一覧。「言語モデルの事前学習は、文章自身から正解を作る自己教師あり学習である」という 1 点だけで十分です。
\end{readingguide}

\chapref{ch:intro}で動かした言語モデルは、「\code{12+35=}」の続きに「\code{47}」と書きました。
足し算の手順を、誰かがプログラムとして書き込んだわけではありません。
モデルは、計算ドリルをたくさん読むうちに、自分で計算を覚えました。

「読むうちに覚える」とは、計算機の中で何が起きることなのでしょうか。
モデルの中には、人間のように「分かった」と感じる部品はありません。
あるのは数の計算だけです。
だとすると、「学習」も何かの計算のはずです。

この章では、その計算が何であるかを、順を追って説明します。
言語モデルの学習は大がかりですが、枠組みは単純で、構成要素は 4 つしかありません。
その枠組みは、「点の集まりに直線を当てはめる」という小さな問題で、すべて確かめられます。
次の章からは、この枠組みはそのままに、モデルを表現力の高いものに取り替えていきます。

\section{ルールを書くか、例から学ぶか}\label{sec:ml-rules-vs-examples}

最初に、ふつうのプログラムと機械学習の違いをはっきりさせます。
この節の問いは、「プログラムを書く代わりに、何をするのか」です。

足し算をする関数なら、\code{def add(a, b): return a + b} と 1 行で書けます。
これは、人が\textbf{ルール}（計算の手順）を知っていて、それをプログラムに書き下したものです。
では、「どうも仕方が」という書きかけの文を受け取り、次に来る言葉を返す関数は書けるでしょうか。
\code{if} 文を並べて書こうとすると、すぐに行き詰まります。
言葉のつながり方のルールは数えきれないほどあり、例外だらけです。
しかも、私たち自身、日本語を話せるのに、そのルールをすべて言葉で説明することはできません。

ルールは書けません。
ところが、\textbf{例}なら大量にあります。
小説を開けば、「どうも仕方が」の次に「ない。」が来る実例が、いくらでも見つかります。
ルールを人が書くのをやめて、例を計算機に見せ、ルールにあたるものを計算機自身に見つけさせるのです。
これが\term{機械学習}{machine learning}です。
\figref{fig:ml-rules-vs-examples}に、二つのやり方の違いを示しました。

\begin{figure}[tbp]
\centering
\begin{tikzpicture}[node distance=6mm and 9mm,
  io/.style={plainbox,minimum width=31mm,minimum height=10mm,font=\footnotesize}]
  % 上段: ふつうのプログラム
  \node[io] (in1) {入力\\\texttt{12, 35}};
  \node[redbox,right=of in1,minimum width=34mm,minimum height=10mm,font=\footnotesize] (rule)
    {\textbf{人が書いたルール}\\\texttt{return a + b}};
  \node[io,right=of rule] (out1) {出力\\\texttt{47}};
  \draw[arr] (in1) -- (rule); \draw[arr] (rule) -- (out1);
  \node[lbl,above=0.5mm of in1.north west,anchor=south west] {\textbf{ふつうのプログラム}：人がルールを書く};
  % 下段: 機械学習
  % 下段の箱は、上段の箱の真下に置く（中心の x 座標を上段とそろえる。縦の矢印をまっすぐにするため）
  \node[orangebox,below=9.5mm of in1,minimum width=31mm,minimum height=13mm,font=\footnotesize] (data)
    {\textbf{入力と正解の例}\\\texttt{(12, 35) → 47}\\\texttt{(5, 65) → 70}\ \ $\cdots$};
  \node[box,minimum width=34mm,minimum height=10mm,font=\footnotesize] (learn) at (rule |- data)
    {\textbf{学習}\\例に合うように調整する};
  \draw[arr] (data) -- (learn);
  \node[io,below=5.5mm of data] (in2) {新しい入力\\\texttt{123, 456}};
  \node[greenbox,minimum width=34mm,minimum height=10mm,font=\footnotesize] (f) at (rule |- in2)
    {\textbf{できあがった関数} $f_\theta$};
  \node[io] (out2) at (out1 |- in2) {予測\\\texttt{579}};
  \draw[arr] (in2) -- (f); \draw[arr] (f) -- (out2);
  \draw[arr] (learn.south) -- node[right,lbl]{関数ができる} (f.north);
  \node[lbl,above=0.5mm of data.north west,anchor=south west] {\textbf{機械学習}：例を見せて、関数を作らせる};
\end{tikzpicture}
\caption{ふつうのプログラム（上段）と機械学習（下段）。機械学習では、人は「入力と正解の例」を用意するだけである。
ルールにあたる関数 $f_\theta$ は学習によって作られ、上段の赤い箱と同じ場所に入って、新しい入力に答える。}
\label{fig:ml-rules-vs-examples}
\end{figure}

「関数を計算機に作らせる」と言っても、何もないところからプログラムが湧いてくるわけではありません。
まず、\textbf{調整できる数値を含んだ関数}を人が用意します。
この数値を変えると、関数のふるまいが変わります。
以下では、この数値を、機械の調整用のつまみにたとえます。
人が用意するのは、\textbf{つまみの付いた関数}です。
次に、例にいちばんよく合うように、つまみを回します。
この「つまみを回す作業」が学習です。

つまみにたとえた数値の正式な名前は、\term{パラメータ}{parameter}です。
パラメータをまとめて $\vtheta$（シータ）と書き、パラメータ $\vtheta$ を持つ関数を $f_\theta$ と書きます。
この「つまみの付いた関数」、つまりパラメータを持つ関数を\term{モデル}{model}と呼びます。
\chapref{ch:intro}の言語モデルは、パラメータを約 260 万個持つ関数でした。

ただし、260 万個のパラメータを、人が手で調整することはできません。
調整のしかたそのものを計算で決めるために、道具があと二つ必要です。

\subsection*{機械学習の 4 要素}

パラメータを自動で調整するには、まず「いまの関数が、例にどれくらい合っていないか」を一つの数で測る必要があります。
この数を計算する関数を\term{損失関数}{loss function}と呼び、その値を\textbf{損失}と呼びます。
損失が小さいほど、関数は例によく合っています。
損失が決まれば、やるべきことは「損失が小さくなるようにパラメータを動かす」ことです。
この作業を\term{最適化}{optimization}と呼びます。

これで、必要なものが四つそろいました。
機械学習は、\textbf{データ}（入力と正解の例）、\textbf{モデル}（パラメータを持つ関数 $f_\theta$）、
\textbf{損失関数}（予測と正解のずれを一つの数で測る）、\textbf{最適化}（損失が小さくなるようにパラメータを動かす）の 4 要素からできています。
\figref{fig:ml-four-elements}は、4 要素のつながりを示しています。
図の矢印を「モデル → 損失関数 → 最適化 → モデル」とたどると、ひと回りして元に戻ります（予測する、損失を計算する、パラメータを更新する）。
設定が適切なら、ひと回りするたびに、損失が少し小さくなる向きへパラメータが動きます。
学習とは、このひと回り（\textbf{学習のループ}）を何百回、何万回と繰り返すことです。

\begin{figure}[tbp]
\centering
\begin{tikzpicture}[node distance=7mm and 13mm,
  el/.style={minimum width=27mm,minimum height=11mm,font=\footnotesize}]
  \node[orangebox,el] (data) {\textbf{1.\ データ}\\入力 $x$ と正解 $y$};
  \node[box,el,right=of data] (model) {\textbf{2.\ モデル} $f_\theta$\\パラメータ $\vtheta$};
  \node[redbox,el,right=of model] (loss) {\textbf{3.\ 損失関数}\\ずれを数にする};
  \node[greenbox,el,below=9mm of model] (opt) {\textbf{4.\ 最適化}\\$\vtheta$ を少し動かす};
  \draw[arr] (data) -- node[above,lbl]{入力 $x$} (model);
  \draw[arr] (model) -- node[above,lbl]{予測} (loss);
  \draw[arr,rounded corners=3pt] (data.north) -- ++(0,6mm) -| node[pos=0.25,above,lbl]{正解 $y$} (loss.north);
  \draw[arr,rounded corners=3pt] (loss.south) |- node[pos=0.75,below,lbl]{損失} (opt.east);
  \draw[arr] (opt) -- node[right,lbl]{$\vtheta$ を更新} (model);
\end{tikzpicture}
\caption{機械学習の 4 要素。データから入力を取り出してモデルに入れ、予測を正解と比べて損失を計算する。
最適化は、損失が小さくなる向きにパラメータ $\vtheta$ を動かす。
「モデル → 損失関数 → 最適化 → モデル」のひと回りを繰り返すことが学習である。}
\label{fig:ml-four-elements}
\end{figure}

この 4 要素は、本書の最後まで変わりません。
この章の最初の例では、データは 100 組の数、モデルは直線、損失関数は「二乗誤差」、最適化は「勾配降下法」です。
\chapref{ch:intro}の言語モデルでは、データが約 250 万トークンの文章に、モデルが約 260 万個のパラメータを持つ関数に替わります。
中身の大きさは桁違いですが、枠組みは同じです。
新しい手法に出会ったら、「4 要素のうち、どれを取り替えたのか」と考えると、位置づけが分かります。

機械学習では、人はルールを書く代わりに 4 要素を用意し、損失が小さくなるようにパラメータを動かす計算を、計算機に繰り返させます。

\section{最初の問題：文字数からトークン数を予測する}\label{sec:ml-supervised}

前節では、4 要素を言葉で説明しました。
この節では、小さな問題を一つ決めて、4 要素のうち「データ」「モデル」「損失関数」を実際に作ります。
問いは、「例によく合う」ということを、どうやって数式にするかです。

\chapref{ch:intro}で見たように、言語モデルは文章を\textbf{トークン}という単位に区切って扱います。
モデルが一度に読める長さも、計算にかかる時間も、文字数ではなくトークン数で決まります。
ところが、ある文章が何トークンになるかは、実際に区切ってみるまで分かりません。
そこで、\textbf{文章の文字数から、トークン数を予測する}という問題を考えます。
「1 文字あたり何トークン」という換算の割合が分かれば、文字数を数えるだけでトークン数を見積もれます。
この割合を、人が決めるのではなく、データから学習させます。

\subsection*{データを作る}

まず例を集めます。
本書の題材である漱石の小説から、長さ 20〜200 字の断片をランダムに切り出します。
それぞれの断片について、文字数と、実際に区切って数えたトークン数を記録します。
区切る道具（トークナイザ）は\chapref{ch:tokenizer}で自作しますが、
ここでは「文章を入れるとトークンの列が返ってくる箱」として使います。
データを作るスクリプトは \file{code/ch03/make_data.py} で、100 件を作りました。
たとえば 1 件目は 118 字の断片で 67 トークン、2 件目は 30 字で 19 トークンです。
100 件すべてを図にすると、\figref{fig:ml-scatter}の左の点になります。
横軸が文字数、縦軸がトークン数です。
点はおおよそ右上がりの直線に沿って並んでいます。
ただし、ぴったり一直線ではありません。
同じくらいの文字数でも、トークン数には上下にばらつきがあります。

\begin{figure}[tbp]
\centering
\begin{tikzpicture}
\begin{axis}[width=0.52\linewidth,height=0.315\linewidth,
  xlabel={文字数},ylabel={トークン数},xmin=0,xmax=210,ymin=0,ymax=150,
  tick label style={font=\scriptsize},label style={font=\footnotesize},
  legend style={font=\scriptsize,at={(0.03,0.97)},anchor=north west,draw=black!30},legend cell align=left]
\addplot[only marks,mark=*,mark size=0.9pt,bookblue] table[col sep=comma,x=chars,y=tokens]
  {code/ch03/out/char_token_train.csv};
\addlegendentry{データ}
\addplot[bookorange,domain=0:210,samples=2,dashed,thick] {0.3*x};
\addlegendentry{$w=0.3$}
\addplot[bookaccent,domain=0:210,samples=2,thick] {0.6*x};
\addlegendentry{$w=0.6$}
\addplot[bookgreen,domain=0:166,samples=2,dashdotted,thick] {0.9*x};
\addlegendentry{$w=0.9$}
\end{axis}
\end{tikzpicture}\hfill
\begin{tikzpicture}
\begin{axis}[width=0.48\linewidth,height=0.315\linewidth,
  xlabel={文字数 $x$（百字）},ylabel={トークン数 $y$（百）},
  xmin=0,xmax=2,ymin=0,ymax=1.2,
  tick label style={font=\scriptsize},label style={font=\footnotesize},
  legend style={font=\scriptsize,at={(0.03,0.97)},anchor=north west,draw=black!30},legend cell align=left]
\addplot[black!60,domain=0:2,samples=2,densely dashed,thick] {0.5*x};
\addlegendentry{$w=0.5$}
\addplot[bookaccent,domain=0:2,samples=2,thick] {0.6*x};
\addlegendentry{$w=0.6$}
\draw[bookorange,line width=1.6pt] (axis cs:0.4,0.20) -- (axis cs:0.4,0.21);
\draw[bookorange,line width=1.6pt] (axis cs:0.8,0.40) -- (axis cs:0.8,0.42);
\draw[bookorange,line width=1.6pt] (axis cs:1.6,0.80) -- (axis cs:1.6,1.01);
\addplot[only marks,mark=*,mark size=1.5pt,bookblue] coordinates {(0.4,0.21) (0.8,0.42) (1.6,1.01)};
\node[font=\scriptsize,text=bookorange,anchor=south east,inner sep=1pt] at (axis cs:1.55,1.0) {$-0.21$};
\node[font=\scriptsize,text=bookorange,anchor=north west] at (axis cs:0.82,0.39) {$-0.02$};
\node[font=\scriptsize,text=bookorange,anchor=north west] at (axis cs:0.42,0.19) {$-0.01$};
\end{axis}
\end{tikzpicture}
\caption{左：文字数とトークン数の関係。青い点がデータ（100 件）で、1 点が 1 つの断片を表す。
3 本の線は、傾き $w$ を変えた直線の候補（切片 $b$ は 0）。
右：3 件のデータ（青い点）と直線。橙の数値が、$w=0.5$ の直線（灰色の破線）の誤差（予測 $-$ 正解）で、橙の縦の線分の長さがその大きさを表す。
3 点とも破線より上にあり（予測が正解より小さい）、誤差は負である。
1・2 件目の線分は短く、点にほぼ隠れている。}
\label{fig:ml-scatter}
\end{figure}

ここで、記号を決めておきます。
モデルに入れる値（文字数）を\textbf{入力}と呼び、$x$ と書きます。
当ててほしい値（トークン数）を\textbf{正解}と呼び、$y$ と書きます。
正解のことを、とくに分類では\term{ラベル}{label}とも呼びます。
データは、入力と正解の組を $N$ 個集めたものです。
$i$ 番目の組を $(x^{(i)}, y^{(i)})$ と書きます。
肩の $(i)$ は「$i$ 番目のデータ」という印で、$i$ 乗という意味ではありません。
データ全体を\term{データセット}{dataset}と呼び、次のように書きます。
\begin{equation}
\calD = \bigl\{(x^{(1)}, y^{(1)}),\ (x^{(2)}, y^{(2)}),\ \dots,\ (x^{(N)}, y^{(N)})\bigr\}
\end{equation}
いまの例では $N=100$ で、$(x^{(1)}, y^{(1)}) = (118, 67)$ です。
このように、入力と正解の組から学ぶやり方を\term{教師あり学習}{supervised learning}と呼びます。
正解が「先生」の役をするからです。

一つだけ、準備をします。
以下の計算では、文字数もトークン数も 100 で割って、「百字」「百トークン」を単位にします。
1 件目は $(x^{(1)}, y^{(1)}) = (1.18, 0.67)$ になります。
数値を 1 前後の大きさにそろえると、計算が安定するからです（理由は\secref{sec:ml-autograd}で説明します）。

\subsection*{モデル：直線}

次にモデルを決めます。
\figref{fig:ml-scatter}の左の点は、ほぼ直線に並んでいました。
そこで、入力 $x$ から予測を作る関数として、直線の式を選びます。
\begin{equation}
f_\theta(x) = wx + b
\label{eq:ml-linear-model}
\end{equation}
$w$ は直線の傾きで、$b$ は切片（$x=0$ のときの値）です。
この二つがパラメータです。まとめて $\vtheta = (w, b)$ と書きます。
$w$ を\term{重み}{weight}、$b$ を\term{バイアス}{bias}と呼びます。
いまの問題では、$w$ は「1 文字あたりのトークン数」を表します。
直線を当てはめてデータを説明するこのモデルを、\term{線形回帰}{linear regression}と呼びます。

$w$ と $b$ を変えると、直線が変わります。
図には、$b=0$ として $w$ を 0.3、0.6、0.9 にした 3 本を描きました。
目で見れば、$w=0.6$ がいちばん良さそうです。
しかし「目で見て良さそう」では、計算機は動けません。
良さを数で表す必要があります。

\subsection*{損失関数：ずれを一つの数にする}

話を追いやすくするため、しばらくデータを 3 件だけに絞ります。
100 件の中から、40 字、80 字、160 字の断片を選びました（\tabref{tab:ml-hand-loss}の左 2 列）。
モデルも一時的に簡単にして、$b=0$ に固定した $f(x)=wx$ を使います。

$w=0.5$ のとき、モデルは 40 字（$x=0.4$）の断片に対して $0.5\times 0.4=0.20$、つまり 20 トークンと予測します。
正解は 21 トークン（$y=0.21$）でした。
予測と正解の差 $0.20-0.21=-0.01$ を\textbf{誤差}と呼びます。
誤差が負なのは、予測が正解より小さいからです。
\figref{fig:ml-scatter}の右では、点と灰色の破線を結ぶ縦の線分の長さが、誤差の大きさ（絶対値）です。
3 件の誤差は $-0.01$、$-0.02$、$-0.21$ です。
3 件目だけが大きく外れています。

3 つの誤差を、一つの数にまとめます。
ここでは、誤差をそれぞれ 2 乗してから平均します。
\tabref{tab:ml-hand-loss}がその計算です。
$w=0.5$ では、$(0.0001+0.0004+0.0441)\div 3 = 0.0446\div 3 = 0.0149$ です。
$w=0.6$ では、同じ計算で $0.0070\div 3 = 0.0023$ になりました。
$w=0.6$ のほうが小さく、目で見た印象と合っています。

\begin{table}[tbp]
\centering
\caption{3 件のデータで、誤差の 2 乗の平均を手で計算する。単位は百字・百トークン。}
\label{tab:ml-hand-loss}
\small
\begin{tabular}{rrrrrrrr}
\toprule
 & & \multicolumn{3}{c}{$w=0.5$} & \multicolumn{3}{c}{$w=0.6$} \\
\cmidrule(lr){3-5}\cmidrule(lr){6-8}
入力 $x$ & 正解 $y$ & 予測 $wx$ & 誤差 & 誤差$^2$ & 予測 $wx$ & 誤差 & 誤差$^2$ \\
\midrule
0.4 & 0.21 & 0.20 & $-0.01$ & 0.0001 & 0.24 & $+0.03$ & 0.0009 \\
0.8 & 0.42 & 0.40 & $-0.02$ & 0.0004 & 0.48 & $+0.06$ & 0.0036 \\
1.6 & 1.01 & 0.80 & $-0.21$ & 0.0441 & 0.96 & $-0.05$ & 0.0025 \\
\midrule
\multicolumn{4}{r}{誤差$^2$ の平均} & 0.0149 & & & 0.0023 \\
\bottomrule
\end{tabular}
\end{table}

この「誤差の 2 乗の平均」を\term{平均二乗誤差}{mean squared error, MSE}と呼びます。
データが $N$ 件のときの式は次のとおりです。
\begin{equation}
\calL(w, b) = \frac{1}{N}\sum_{i=1}^{N}\bigl(f_\theta(x^{(i)}) - y^{(i)}\bigr)^2
            = \frac{1}{N}\sum_{i=1}^{N}\bigl(wx^{(i)} + b - y^{(i)}\bigr)^2
\label{eq:ml-mse}
\end{equation}
$\calL$ は損失（loss）の頭文字です。
この式で大事なのは、$\calL$ が\textbf{パラメータ $w, b$ の関数}だという点です。
データ $x^{(i)}, y^{(i)}$ は、集めた時点で決まった数です。
動かせるのは $w$ と $b$ だけです。
$w$ と $b$ を決めると直線が決まり、直線が決まると損失が一つ決まります。

\begin{quiz}
誤差を 2 乗せずに、そのまま平均してはいけないのでしょうか。
2 乗の代わりに絶対値を使うのはどうでしょうか。
\end{quiz}

そのまま平均すると、正の誤差と負の誤差が打ち消し合います。
誤差が $+0.5$ と $-0.5$ の 2 件は、平均が 0 になり、「完璧」と判定されてしまいます。
2 乗でも絶対値でも、打ち消し合いは防げます。実際、絶対値を使う流儀もあります。
ただし 2 乗には、(1) なめらかで微分しやすい、(2) 大きな外れを特に重く数える、という性質があります。
\tabref{tab:ml-hand-loss}の $w=0.5$ では、3 件目の誤差は 1 件目の 21 倍ですが、2 乗すると 441 倍になっています。
2 乗を選ぶことには、じつはもっと深い理由があります（\secref{sec:ml-mle}）。

\subsection*{一般の形：損失の平均を最小にする}

ここまでの話を、問題によらない形で書いておきます。
1 件のデータについて、予測 $f_\theta(x)$ と正解 $y$ のずれを測る関数を $\ell$（エル）と書きます。
いまの例では $\ell(f_\theta(x), y) = (f_\theta(x) - y)^2$ です。
全体の損失は、1 件ごとの損失の平均です。

\begin{definition}[教師あり学習で最小にする量]\label{def:ml-erm}
データセット $\calD = \{(x^{(i)}, y^{(i)})\}_{i=1}^{N}$ と、モデル $f_\theta$ と、1 件ごとの損失 $\ell$ が与えられたとする。
学習とは、損失の平均
\begin{equation}
\calL(\vtheta) = \frac{1}{N}\sum_{i=1}^{N} \ell\bigl(f_\theta(x^{(i)}),\, y^{(i)}\bigr)
\label{eq:ml-erm}
\end{equation}
がなるべく小さくなるパラメータ $\vtheta$ を求めることである。
\end{definition}

この方針には、\term{経験リスク最小化}{empirical risk minimization}という名前が付いています。
「リスク」は、まだ見ていないデータも含めた、損失の平均です。
これは計算できないので、手元の $N$ 件での平均で代用します
（「経験」はこの代用を指します。どこまで通用するかは、\secref{sec:ml-generalization}で調べます）。
本書に出てくる学習は、言語モデルの事前学習も含めて、ほとんどがこの形をしています。
変わるのは、$f_\theta$ と $\ell$ の中身だけです。

この節では、4 要素のうち三つを作りました。
データは（文字数, トークン数）の組、モデルは直線 $wx+b$、損失関数は平均二乗誤差です。
「例によく合う」とは、「損失 $\calL$ が小さい」ということです。
残る一つは、損失を小さくする $w$ と $b$ を探す最適化です。

\section{勾配降下法：損失の坂を下る}\label{sec:ml-gd}

前節で、直線の良さを測る数（損失）ができました。
この節の問いは、「損失が最小になるパラメータを、計算機にどう探させるか」です。
答えは\term{勾配降下法}{gradient descent}です。
本書の最後まで、学習はすべてこの方法とその改良版で行います。

\subsection*{パラメータが 1 個のとき}

まず、前節の 3 件のデータと、$b=0$ に固定したモデル $f(x)=wx$ で考えます。
\eqref{eq:ml-mse}に 3 件の値を入れて展開すると、損失は次の式になります。
\begin{align}
\calL(w) &= \tfrac{1}{3}\bigl\{(0.4w-0.21)^2 + (0.8w-0.42)^2 + (1.6w-1.01)^2\bigr\} \notag\\
         &= 1.12\,w^2 - 1.3573\,w + 0.4135
\label{eq:ml-loss-1d}
\end{align}
2 行目の係数の出どころを確かめておきます。
1 件目の $(0.4w-0.21)^2$ を展開すると、$0.16\,w^2 - 2\times0.4\times0.21\,w + 0.21^2 = 0.16\,w^2 - 0.168\,w + 0.0441$ です。
残りの 2 件も同じように展開し、$w^2$ の項、$w$ の項、定数の項ごとに集めて 3 で割ります。
\begin{align*}
w^2\ \text{の係数} &= (0.16 + 0.64 + 2.56)\div 3 = 3.36\div 3 = 1.12\\
w\ \text{の係数} &= -(0.168 + 0.672 + 3.232)\div 3 = -4.072\div 3 = -1.3573\\
\text{定数の項} &= (0.0441 + 0.1764 + 1.0201)\div 3 = 1.2406\div 3 = 0.4135
\end{align*}
$w$ の 2 次式なので、グラフは\figref{fig:ml-loss-1d}のような放物線です。
$w^2$ の係数が正なので、谷の形をしています。
求めたいのは、損失がいちばん小さくなる $w$ です。
グラフでいちばん低いこの点を、以下では\textbf{谷底}と呼びます。

\begin{figure}[tbp]
\centering
\begin{tikzpicture}
\begin{axis}[width=0.72\linewidth,height=0.29\linewidth,
  xlabel={重み $w$},ylabel={損失 $\calL(w)$},xmin=-0.15,xmax=1.35,ymin=-0.06,ymax=0.52]
\addplot[bookblue,thick,domain=-0.15:1.35,samples=80] {1.12*x*x-1.3573*x+0.4135};
\addplot[only marks,mark=*,mark size=1.6pt,bookblue] coordinates
  {(0,0.4135) (0.3,0.1071) (0.6,0.0023) (0.9,0.0991) (1.2,0.3975)};
\addplot[bookaccent,very thick,domain=-0.1:0.2,samples=2] {0.4135-1.3573*x};
\addplot[bookgreen,very thick,domain=0.7:1.1,samples=2] {0.0991+0.6587*(x-0.9)};
\node[font=\footnotesize,text=bookaccent,anchor=west] at (axis cs:0.07,0.42) {傾き $-1.36$：$w$ を増やせば下がる};
\node[font=\footnotesize,text=bookgreen,anchor=west] at (axis cs:0.2,0.29) {傾き $+0.66$：$w$ を減らせば下がる};
\draw[bookgreen,thin] (axis cs:0.86,0.25) -- (axis cs:0.9,0.12);
\draw[black,thin] (axis cs:0.606,0.075) -- (axis cs:0.606,0.018);
\node[font=\footnotesize,anchor=south] at (axis cs:0.606,0.07) {谷底 $w=0.606$};
\end{axis}
\end{tikzpicture}
\caption{3 件のデータに対する損失 $\calL(w)$。青い点は、$w$ を 0.3 刻みで 5 通り試したときの損失。
太い線分は、$w=0$（赤）と $w=0.9$（緑）での接線で、その場所での傾き（微分）を表す。
傾きが負なら右へ、正なら左へ動けば、損失は下がる。}
\label{fig:ml-loss-1d}
\end{figure}

いちばん素朴な探し方は、$w$ を細かく刻んで全部試すことです。
\figref{fig:ml-loss-1d}の青い点は、0.3 刻みで 5 通りを試した結果で、$w=0.6$ の損失がいちばん小さくなっています。
しかし、この方法はパラメータが増えると使えなくなります。
1 個あたり 100 通りを試すなら、2 個で $100^2 = 1$ 万通り、260 万個では $100^{2600000}$ 通りで、とても終わりません。

そこで、\textbf{いまいる場所の傾き}だけを手がかりにします。
この方法を、霧の深い山で谷底を目指す場面にたとえます（山の高さが損失、人の位置がいまの $w$、足元の傾きが微分）。
霧で遠くは見えませんが、足元がどちらに下っているかは分かります。
同じように、損失のグラフ全体を計算しなくても、いまの $w$ での傾きは計算できます。
下っている向きに一歩進み、また足元を確かめる。
これを繰り返せば、谷底に近づけます。
節の題「損失の坂を下る」は、このたとえです（山と違い、パラメータが増えると位置は数の組です）。

グラフの傾きとは、微分のことでした（\chapref{ch:math}）。
\eqref{eq:ml-loss-1d}を $w$ で微分します。
$w^2$ の微分は $2w$、$w$ の微分は $1$、定数の微分は $0$ です。
したがって、$\dd{\calL}{w} = 1.12\times 2w - 1.3573 = 2.24\,w - 1.3573$ です。
$w=0$ での傾きは $-1.3573$ です。
傾きが負なので、$w$ を増やすと損失が減ります。
$w=0.9$ での傾きは $2.24\times 0.9 - 1.3573 = +0.6587$ です。
こんどは正なので、$w$ を減らすと損失が減ります。
どちらの場合も、\textbf{傾きの符号と逆の向き}に少しだけ動けば、損失は下がります。
これを一つの式にまとめたものが、次の更新の規則です。
\begin{equation}
w \leftarrow w - \eta\,\dd{\calL}{w}
\label{eq:ml-gd-1d}
\end{equation}
矢印 $\leftarrow$ は、「右辺を計算して、その結果で $w$ を置き換える」という意味です。
$\eta$（イータ）は、1 回の更新の幅を決める正の数で、\term{学習率}{learning rate}と呼びます。
1 回の更新を 1 \term{ステップ}{step}と数えます。

$\eta=0.4$ として、$w=0$ から手で計算してみます。
1 ステップ目は $w = 0 - 0.4\times(-1.3573) = 0.5429$ です。
この場所での傾きは $2.24\times 0.5429 - 1.3573 = -0.1412$ です。
2 ステップ目は $w = 0.5429 - 0.4\times(-0.1412) = 0.5994$ となります。
同じ計算をコードで 5 ステップ続けました（\file{code/ch03/hand_example.py}）。
$w$ は $0.5429$、$0.5994$、$0.6053$、$0.6059$ と進み、$0.606$ のあたりで動かなくなりました。
この $0.606$ が、\figref{fig:ml-loss-1d}の谷底です。
谷底では傾きが 0 なので、$2.24\,w - 1.3573 = 0$ を解くと、$w = 1.3573\div 2.24 = 0.606$ と確かめられます。
この間、傾きは $-1.357$、$-0.141$、$-0.015$、$-0.002$ と、谷底に近づくほど 0 に近づいています。
更新の幅は「学習率 $\times$ 傾き」なので、谷底に近づくと更新の幅も自然に小さくなり、谷底で $w$ は止まります。

学習率 $\eta$ は、人が決める数です。
\figref{fig:ml-lr-1d}は、$\eta$ を変えて同じ計算を 5 ステップ行った様子です。
$\eta=0.05$（左）では、1 回の更新の幅が小さすぎて、5 ステップではまだ谷底まで半分も進んでいません。
$\eta=0.4$（中央）では、2 ステップでほぼ谷底に着きます。
$\eta=1.0$（右）では、谷底を飛び越えて放物線の反対側に移り、しかも元より損失の大きい場所に着いています。
1 ステップ目は $w = 0 - 1.0\times(-1.3573) = 1.3573$ で、損失は $0.4135$ から $0.635$ に増えました。
これを繰り返すので、損失は増え続けます。これを\textbf{発散}と呼びます。

\begin{figure}[tbp]
\centering
\pgfplotsset{mlsmall/.style={width=0.385\linewidth,height=0.29\linewidth,
  title style={font=\footnotesize,yshift=-1.5mm},tick label style={font=\scriptsize},
  label style={font=\scriptsize},xlabel={$w$}}}
\begin{tikzpicture}
\begin{axis}[mlsmall,title={$\eta=0.05$：遅い},ylabel={損失},xmin=-0.2,xmax=1.4,ymin=0,ymax=0.6]
\addplot[black!45,domain=-0.2:1.4,samples=50] {1.12*x*x-1.3573*x+0.4135};
\addplot[bookaccent,thick,mark=*,mark size=1.3pt] table[col sep=comma,x=w,y=loss]{code/ch03/out/hand_gd_eta0.05.csv};
\node[font=\scriptsize,anchor=south west,inner sep=1pt] at (axis cs:0.04,0.43) {出発点};
\end{axis}
\end{tikzpicture}\hfill
\begin{tikzpicture}
\begin{axis}[mlsmall,title={$\eta=0.4$：ちょうどよい},xmin=-0.2,xmax=1.4,ymin=0,ymax=0.6]
\addplot[black!45,domain=-0.2:1.4,samples=50] {1.12*x*x-1.3573*x+0.4135};
\addplot[bookaccent,thick,mark=*,mark size=1.3pt] table[col sep=comma,x=w,y=loss]{code/ch03/out/hand_gd_eta0.4.csv};
\node[font=\scriptsize,anchor=south west,inner sep=1pt] at (axis cs:0.04,0.43) {出発点};
\end{axis}
\end{tikzpicture}\hfill
\begin{tikzpicture}
\begin{axis}[mlsmall,title={$\eta=1.0$：発散},xmin=-1.2,xmax=2.8,ymin=0,ymax=4]
\addplot[black!45,domain=-1.2:2.8,samples=50] {1.12*x*x-1.3573*x+0.4135};
\addplot[bookaccent,thick,mark=*,mark size=1.3pt] table[col sep=comma,x=w,y=loss]{code/ch03/out/hand_gd_eta1.0.csv};
% ステップの番号（放物線の外側に置く）
\node[font=\scriptsize,anchor=east,inner sep=1pt] at (axis cs:-0.12,0.42) {0};
\node[font=\scriptsize,anchor=north west,inner sep=1pt] at (axis cs:1.43,0.58) {1};
\node[font=\scriptsize,anchor=east,inner sep=2pt] at (axis cs:-0.3258,0.9746) {2};
\node[font=\scriptsize,anchor=west,inner sep=2pt] at (axis cs:1.7613,1.4972) {3};
\node[font=\scriptsize,anchor=east,inner sep=2pt] at (axis cs:-0.8266,2.3009) {4};
\node[font=\scriptsize,anchor=west,inner sep=2pt] at (axis cs:2.3824,3.5367) {5};
\node[font=\scriptsize,anchor=south,inner sep=1pt] at (axis cs:0.25,3.2) {0 が出発点};
\end{axis}
\end{tikzpicture}
\caption{学習率 $\eta$ による違い。灰色の曲線が損失、赤い点が、出発点 $w=0$ とそこからの 5 ステップ（計 6 点）で、赤い線が進んだ順を表す。
中央は、2 ステップ目以降の点が谷底でほぼ重なっている。右は、点の横の数字がステップの番号。
右の図だけ、軸の範囲が広いことに注意。}
\label{fig:ml-lr-1d}
\end{figure}

学習率は、大きすぎても小さすぎてもいけません。
どこからが大きすぎるかは、谷の急さで決まり、問題ごとに違います。
この例では、傾きが $2.24\,(w - 0.606)$ と書けるので、1 ステップで、谷底からの距離は $(1 - 2.24\,\eta)$ 倍になります。
この倍率の絶対値が 1 を超える $\eta > 0.89$ で発散します。
ふつうは、この境目を計算では求められません。
そこで実際の学習では、いくつかの値を試して、損失がいちばん速く安定して下がるものを選びます。

\subsection*{パラメータが 2 個のとき：勾配}

100 件のデータと、$w$ と $b$ の両方を持つモデルに戻ります。
パラメータが 2 個になると、損失は、$w$ と $b$ の二つを決めると値が一つ決まる関数になります。
損失の値を高さと考えると、$(w, b)$ の平面の上に、高い所と低い所のある地形ができます。
\figref{fig:ml-landscape}の左は、この地形を真上から見て、地図と同じように等高線で描いたものです。
楕円が等高線で、同じ損失になる $(w, b)$ をつないだものです。
内側の楕円ほど損失が小さく、中心の $\times$ 印が谷底です。

考え方は 1 個のときと同じです。
違うのは、傾きが $w$ 方向と $b$ 方向の二つあることです。
$b$ を止めて $w$ だけを動かしたときの傾きが偏微分 $\pd{\calL}{w}$、
$w$ を止めて $b$ だけを動かしたときの傾きが $\pd{\calL}{b}$ でした（\chapref{ch:math}）。
この二つを\eqref{eq:ml-mse}から求めます。
式を見やすくするため、$i$ 番目のデータの誤差を $e^{(i)} = wx^{(i)} + b - y^{(i)}$ と置きます。
すると損失は $\calL = \frac{1}{N}\sum_i (e^{(i)})^2$ です。

まず $w$ で微分します。
和の微分は、項ごとの微分の和です。
各項 $(e^{(i)})^2$ は、「$e^{(i)}$ を 2 乗する」と「$w$ から $e^{(i)}$ を作る」の合成です。
合成関数の微分により、$2e^{(i)}$ に、$e^{(i)}$ を $w$ で微分したものを掛けます。
$e^{(i)} = wx^{(i)} + b - y^{(i)}$ を $w$ で微分すると $x^{(i)}$ です。
$b$ で微分すると 1 です。
したがって、次の二つの式が得られます。
\begin{equation}
\pd{\calL}{w} = \frac{2}{N}\sum_{i=1}^{N} e^{(i)}\,x^{(i)},
\qquad
\pd{\calL}{b} = \frac{2}{N}\sum_{i=1}^{N} e^{(i)}
\label{eq:ml-linreg-grad}
\end{equation}
言葉で言えば、$\pd{\calL}{w}$ は「誤差 $\times$ 入力」の平均の 2 倍、$\pd{\calL}{b}$ は「誤差」の平均の 2 倍です。
この「誤差 $\times$ 入力」という形は、このあと何度も現れます。

二つの偏微分を並べたベクトルが、\term{勾配}{gradient} $\nabla_{\vtheta}\calL$ でした。
勾配は、損失が最も急に\textbf{増える}向きを指します（\chapref{ch:math}）。
損失を減らすには、その逆向きに進みます。

\begin{definition}[勾配降下法]\label{def:ml-gd}
損失 $\calL(\vtheta)$ を小さくするために、パラメータを適当な初期値から始めて、次の更新を繰り返す。
\begin{equation}
\vtheta \leftarrow \vtheta - \eta\,\nabla_{\vtheta}\calL(\vtheta)
\label{eq:ml-gd}
\end{equation}
$\eta > 0$ は学習率である。
線形回帰では $\vtheta=(w,b)$ なので、
$w \leftarrow w - \eta\,\pd{\calL}{w}$ と $b \leftarrow b - \eta\,\pd{\calL}{b}$ を同時に行う。
\end{definition}

パラメータが 260 万個になっても、式は\eqref{eq:ml-gd}のままです。
勾配が 260 万個の成分を持つベクトルになるだけです。
勾配降下法の 1 ステップに要るのは、勾配一つです。
その全成分は、損失を 1 回計算する手間の数倍で求まります（\chapref{ch:nn}の誤差逆伝播法）。

ここまでを NumPy で実装します。
100 件のデータを読み込み、100 で割った配列を \code{x}、\code{y} とします。
どちらも形は $(N,)$ です。
まず、モデルと損失関数と勾配です。

\snippet{線形回帰のモデル・損失・手で導いた勾配}{code:ml-linreg-model}{ch03/linreg_numpy.py}{model-loss}

3 行目が\eqref{eq:ml-linear-model}のモデルです。
\code{x} は 100 件分の配列なので、\code{w * x + b} は 100 件の予測をまとめて計算します（\chapref{ch:pytorch}のブロードキャスト）。
8 行目が\eqref{eq:ml-mse}です。予測から正解を引き、2 乗し、\code{np.mean} で平均します。
13 行目の \code{err} が、100 件の誤差 $e^{(i)}$ を並べた配列です。
14 行目と 15 行目が、\eqref{eq:ml-linreg-grad}の二つの式です。
\code{err * x} は、誤差と入力を 1 件ずつ掛け合わせます。

手で導いた式には、間違いが入りこみがちです。
そこで、パラメータを少し動かして損失の変化を直接測る数値微分（\chapref{ch:math}）で検算しました。
点 $(w,b)=(0.3, 0.2)$ では、どちらの方法でも勾配は $(-0.503841, -0.279140)$ で、小数第 6 位まで一致しました。
次が、勾配降下法の本体です。

\snippet{勾配降下法}{code:ml-gd}{ch03/linreg_numpy.py}{gd}

4 行目からのループが、\eqref{eq:ml-gd}の繰り返しです。
5 行目で、いまの $(w,b)$ での勾配を求めます。
6〜7 行目で、勾配の逆向きに $\eta$ 倍だけ動かします。
3 行目と 8 行目は、あとで図を描くために途中経過を記録しているだけです。
初期値 $(w,b)=(0,0)$、学習率 $\eta=0.1$ で 1000 ステップ動かした結果が次です。

\outputfile{ch03/out/linreg_run.txt}

損失は、最初の 0.587 から、2 ステップで 0.061、100 ステップで 0.0113 まで下がりました。
300 ステップ以降は、$w=0.59$、$b=0.008$、損失 $0.01121$ のあたりからほとんど動きません。
1 ステップ目を手で確かめておきましょう。
$(w,b)=(0,0)$ では予測がすべて 0 なので、誤差は $e^{(i)}=-y^{(i)}$ です。
これを\eqref{eq:ml-linreg-grad}に入れると、$\pd{\calL}{w}$ は「$x^{(i)}y^{(i)}$ の平均」の $-2$ 倍、$\pd{\calL}{b}$ は「$y^{(i)}$ の平均」の $-2$ 倍です。
100 件のデータでは、$x^{(i)}y^{(i)}$ の平均が $0.9657$、$y^{(i)}$ の平均が $0.6816$ です。
したがって、勾配は $(-2\times0.9657,\ -2\times0.6816) = (-1.931, -1.363)$ になります。
更新は $w = 0 - 0.1\times(-1.931) = 0.1931$、$b = 0 - 0.1\times(-1.363) = 0.1363$ で、出力の 2 行目と一致します。

\begin{figure}[tbp]
\centering
\begin{tikzpicture}
\begin{axis}[width=0.5\linewidth,height=0.315\linewidth,xlabel={重み $w$},ylabel={バイアス $b$},
  xmin=-0.25,xmax=1.25,ymin=-0.75,ymax=0.75,
  tick label style={font=\scriptsize},label style={font=\footnotesize}]
\addplot[bookblue!55,thin,empty line=jump] table[col sep=comma,x=w,y=b]{code/ch03/out/loss_contours.csv};
\addplot[bookaccent,thick] table[col sep=comma,x=w,y=b]{code/ch03/out/linreg_path.csv};
\addplot[bookaccent,only marks,mark=*,mark size=1.2pt,restrict expr to domain={\thisrow{step}}{0:10}]
  table[col sep=comma,x=w,y=b]{code/ch03/out/linreg_path.csv};
\addplot[only marks,mark=x,mark size=3pt,thick,black] coordinates {(0.5906,0.0083)};
\node[font=\scriptsize,anchor=north east] at (axis cs:0.02,-0.02) {出発点};
\node[font=\scriptsize,anchor=south west] at (axis cs:0.62,0.03) {谷底};
\coordinate (ml-c01) at (axis cs:0.085,0.672);
\end{axis}
\draw[bookblue,thin] (ml-c01) -- ++(0,2.6mm)
  node[above,font=\scriptsize,text=bookblue,inner sep=1pt] {損失 0.1};
\end{tikzpicture}\hfill
\begin{tikzpicture}
\begin{axis}[width=0.46\linewidth,height=0.315\linewidth,xlabel={文字数},ylabel={トークン数},
  xmin=0,xmax=210,ymin=0,ymax=150,clip=false,
  tick label style={font=\scriptsize},label style={font=\footnotesize}]
\addplot[only marks,mark=*,mark size=0.8pt,bookblue!60] table[col sep=comma,x=chars,y=tokens]
  {code/ch03/out/char_token_train.csv};
\addplot[black!50,domain=0:210,samples=2] {0.1931*x+13.63} node[pos=1,anchor=west,font=\scriptsize] {1};
\addplot[black!50,domain=0:210,samples=2] {0.2926*x+20.13} node[pos=1,anchor=west,font=\scriptsize] {2};
\addplot[black!50,domain=0:210,samples=2] {0.4252*x+22.37} node[pos=1,anchor=west,font=\scriptsize,yshift=-1.5pt] {10};
\addplot[bookaccent,very thick,domain=0:210,samples=2] {0.5906*x+0.83}
  node[pos=1,anchor=west,font=\scriptsize,text=bookaccent,yshift=1.5pt] {1000};
\end{axis}
\end{tikzpicture}
\caption{勾配降下法の二つの見方。左：損失の等高線（青）と、パラメータ $(w,b)$ が動いた跡（赤）。
丸印は、出発点と最初の 10 ステップ（計 11 個。4 ステップ目より後の点は、間隔が丸印の大きさより狭く、重なって 1 つの塊に見える）で、$\times$ 印が谷底。
等高線は、内側から損失 0.02、0.05、0.1、0.2、0.4、0.8、1.6（0.2 から外の線は、枠で切れている）。
右：同じ学習を、データの上の直線で見たもの。右端の数字はステップ数で、赤い線が学習後（1000 ステップ後）。}
\label{fig:ml-landscape}
\end{figure}

\figref{fig:ml-landscape}は、この学習を二つの見方で描きました。
左の図では、パラメータは最初の数ステップで等高線を何本も横切り、損失の小さい側へ大きく動きます。
そのあとは向きを変え、細長い谷に沿ってゆっくり谷底へ向かいます。
右の図で見ると、意味が分かります。
最初の数ステップで、直線は点の集まりの真ん中まで持ち上がります。
そのあとは、真ん中を通ったまま、傾きを少しずつ直しています。
点の集まりの真ん中さえ通っていれば、傾きが多少違っても損失はあまり変わりません。
谷が細長いのは、このためです。

学習率を $\eta=0.42$ まで上げると、100 ステップ後の損失は $2.7\times 10^{7}$ になり、発散しました。
パラメータが 2 個でも、\figref{fig:ml-lr-1d}と同じことが起きます。
$\eta$ を細かく変えて試すと、収束と発散の境目は約 $0.40$ でした。
3 件のときの境目（約 $0.89$）と違うのは、データとパラメータの数が変わって、谷の急さが変わったからです。

% [figpass] 案 5
学習率を細かく変えた結果を、\figref{fig:ml-sweep-lr}にまとめました。
\file{code/ch03/sweep_lr.py} で、出発点（0）とステップ数を固定し、学習率だけを、3 件のデータでは 0.02 刻みで 55 通り、100 件のデータでは 0.01 刻みで 45 通りに変えて走らせたものです。
どちらも、学習率が小さすぎる左端では損失が下がりきらず、中ほどに損失の低い広い範囲があり、境目を越えると損失が急に跳ね上がります。
灰色の点線で示した境目（0.893 と 0.401）は、損失の式の 2 次の係数から計算した値で、跳ね上がる位置と一致します。
左の図で、5 ステップ後の損失（赤）が境目の手前から上がり始めるのは、谷底をはさんで左右に行き来しながら近づくので、5 ステップでは足りないためです。
50 ステップ後（青の破線）では、学習率 0.86 までは谷底に着いています。

\begin{figure}[tbp]
\centering
\pgfplotsset{mlsweeplr/.style={height=0.33\linewidth,
  ymin=3e-4,ymax=20,xlabel={学習率 $\eta$},
  tick label style={font=\scriptsize},label style={font=\footnotesize},
  title style={font=\footnotesize,yshift=-1.5mm},
  legend style={font=\scriptsize,at={(0.03,0.97)},anchor=north west,draw=black!30},legend cell align=left}}
\begin{tikzpicture}
\begin{semilogyaxis}[mlsweeplr,width=0.53\linewidth,title={3 件のデータ（パラメータ $w$ だけ）},
  ylabel={損失},xmin=0,xmax=1.1]
\addplot[bookaccent,thick,mark=*,mark size=0.8pt] table[col sep=comma,x=eta,y=loss5]{code/ch03/out/sweep_lr_1d.csv};
\addplot[bookblue,thick,densely dashed] table[col sep=comma,x=eta,y=loss50]{code/ch03/out/sweep_lr_1d.csv};
\legend{5 ステップ後,50 ステップ後}
\draw[black!50,thick,densely dotted] (axis cs:0.893,3e-4) -- (axis cs:0.893,20);
\node[font=\scriptsize,anchor=south east,inner sep=1pt] at (axis cs:0.88,3.5e-4) {境目 0.893};
\end{semilogyaxis}
\end{tikzpicture}\hfill
\begin{tikzpicture}
\begin{semilogyaxis}[mlsweeplr,width=0.49\linewidth,title={100 件のデータ（$w$ と $b$）},
  xmin=0,xmax=0.45,xtick={0,0.1,0.2,0.3,0.4}]
\addplot[bookaccent,thick,mark=*,mark size=0.8pt] table[col sep=comma,x=eta,y=loss100]{code/ch03/out/sweep_lr_2d.csv};
\legend{100 ステップ後}
\draw[black!50,thick,densely dotted] (axis cs:0.401,3e-4) -- (axis cs:0.401,20);
\node[font=\scriptsize,anchor=south east,inner sep=1pt] at (axis cs:0.395,3.5e-4) {境目 0.401};
\end{semilogyaxis}
\end{tikzpicture}
\caption{学習率 $\eta$ だけを変えたときの、決まった回数の更新のあとの損失（縦軸は対数目盛り）。出発点は 0 に固定した。
左は\figref{fig:ml-lr-1d}と同じ 3 件のデータ（0.02 刻み）、右は 100 件のデータ（0.01 刻み）。
小さすぎると下がりきらず、境目（灰色の点線。損失の式から計算した値）を越えると発散して、損失が枠の上へ飛び出す。
\figref{fig:ml-lr-1d}の 3 枚のうち、$\eta=0.4$ と $1.0$ は左の図の点にあたり、$\eta=0.05$ は $0.04$ と $0.06$ の 2 点の間にあたる（0.02 刻みなので、$0.05$ の点はない）。}
\label{fig:ml-sweep-lr}
\end{figure}
% [/figpass]

得られた直線は $y = 0.5906\,x + 0.0083$ です。
$w = 0.59$ は、「1 文字あたり約 0.59 トークン」を意味します。
$b=0.0083$ は百トークン単位なので、0.8 トークンほどで、ほぼ 0 です。
つまり、「100 字の文章は約 60 トークン」と見積もればよいことが、データから分かりました。
損失は誤差の 2 乗の平均なので、平方根を取ると、誤差と同じ単位に戻ります。
$\sqrt{0.01121} = 0.106$ 百トークン、つまり約 10 トークンです。
これが、この直線の典型的な誤差の大きさです。

じつは線形回帰に限っては、勾配が 0 になる条件を連立方程式として解き、最小点を一度で求める公式があります（\textbf{正規方程式}）。
この公式でも、同じ $w=0.5906$、$b=0.0083$ になります。
それでも勾配降下法を使うのは、次章以降のモデルにはそのような公式がなく、勾配なら計算できるからです。

この節では、損失を少しずつ小さくしていく方法（損失の坂を下る方法）を作りました。
勾配を求め、その逆向きに学習率の分だけ動く。これを繰り返すだけです。
4 要素の最後の一つ、最適化がそろいました。

\section{勾配を自動で求める}\label{sec:ml-autograd}

前節では、勾配の式を手で導きました。
モデルを変えるたびに、式を導き直さなければならないのでしょうか。
この節の答えは「いいえ」です。
\chapref{ch:pytorch}で学んだ PyTorch の自動微分に任せられます。

\code{requires_grad=True} を付けたテンソルから損失を計算して \code{backward()} を呼ぶと、そのテンソルの \code{.grad} に勾配が入るのでした。
そこで、$w$ と $b$ を、どちらも \code{torch.tensor(0.0, requires_grad=True)} で作ります。
この $w$ と $b$ から損失を計算して \code{backward()} を呼ぶと、勾配は $(-1.931374,\ -1.363200)$ でした。
\eqref{eq:ml-linreg-grad}で計算した値と、小数第 6 位まで一致します。
勾配の式は、コードのどこにも書いていません。
学習のループは次のとおりです。\code{x} と \code{y} は、データをテンソルに直したものです。

\snippet{自動微分を使った学習ループ}{code:ml-torch-loop}{ch03/linreg_torch.py}{loop}

ループは 4 手順です。
(1) 予測して損失を計算します（3 行目）。
(2) \code{backward()} で勾配を計算します（7 行目）。
(3) 勾配の逆向きにパラメータを動かします（8〜10 行目）。
更新は微分の対象ではないので、\code{torch.no_grad()} で囲みます。
(4) 勾配を 0 に戻します（11〜12 行目）。
\code{backward()} は勾配を前の値に足し込むからでした（\chapref{ch:pytorch}）。
4〜6 行目は途中経過の表示です（繰り返しが 1001 回なのは、1000 回更新したあとの値も表示するためです）。

実行すると、1000 ステップ後に $w=0.5906$、$b=0.0083$、損失 $0.01121$ となりました。
NumPy 版とまったく同じ数値です。
大事なのは、このループが\textbf{モデルの中身を知らない}ことです。
3 行目の損失の計算を別のモデルのものに取り替えても、(2)〜(4) は 1 文字も変わりません。
本書の学習コードは、最後の章まで、すべてこの 4 手順の繰り返しです。
次章では、(3) の更新と (4) の後始末を PyTorch の部品に任せて、ループを 5 行に書き直します。
行の数え方が変わるだけで、していることは、この 4 手順と同じです。

\subsection*{足し算を、例だけから学ぶ}

章の冒頭の問いに戻りましょう。
足し算のルールを書かずに、例だけから足し算を学べるでしょうか。
いまある道具で試せます。
計算ドリルから足し算の問題を 1,000 題取り出し、二つの数を入力、答えを正解とします。
「\code{12+35=47}」という行からは、入力 $(12, 35)$ と正解 $47$ ができます。

入力が二つになったので、モデルを少し広げます。
入力を $x_1, x_2$ とし、それぞれに重みを掛けて足します。
\begin{equation}
f_\theta(\vx) = w_1 x_1 + w_2 x_2 + b = \vw^\T\vx + b
\label{eq:ml-linear-multi}
\end{equation}
$\vx=(x_1,x_2)^\T$ は入力を並べたベクトル、$\vw=(w_1,w_2)^\T$ は重みを並べたベクトルです。
$\vw^\T\vx$ は二つの内積です（\chapref{ch:math}）。
入力が何個に増えても、この形で書けます。
学習によって $w_1=1$、$w_2=1$、$b=0$ になれば、このモデルは足し算になります。

データは 1000 で割り、入力 \code{X}（形 $(N,2)$）と正解 \code{y}（形 $(N,)$）にします（\file{code/ch03/learn_addition.py}）。
学習のループは\coderef{code:ml-torch-loop}と同じです。
変えるのは、パラメータの用意と、損失の計算の 2 か所だけです。

\begin{pycode}
w = torch.zeros(2, requires_grad=True)  # 重みは入力ごとに 1 つ: (w1, w2)
y_hat = X @ w + b  # 全データの予測。形 (N,)
loss = ((y_hat - y) ** 2).mean()
\end{pycode}

1 行目はループの前に置くパラメータの用意で、2〜3 行目がループの中の損失の計算です。
バイアス \code{b} の用意は、これまでと同じなので省きました。
\code{X @ w} は、1,000 件それぞれについて内積 $\vw^\T\vx$ を計算します。
学習率 $0.5$ で 1000 ステップ学習した結果が次です
（前節で見たとおり、ちょうどよい学習率は問題ごとに違うので、線形回帰のときとは別の値です）。

\outputfile{ch03/out/learn_addition.txt}

1 行目は、データの最初の 3 件です。
損失の \code{2.73e-01} は、$2.73\times10^{-1}$ という数の書き方です。
重みは $(0,0)$ から始まり、10 ステップで $(0.5510, 0.5279)$、100 ステップで $(0.9993, 0.9991)$ になりました。
バイアスは、10 ステップ目にいったん $0.1764$ まで増えますが、100 ステップ目には $0.0003$ に戻っています。
最初は重みが小さく、予測が正解より小さいので、まずバイアスが増えて予測全体を持ち上げます。
重みが 1 に近づくと、その分が要らなくなり、バイアスは 0 に戻ります（\figref{fig:ml-landscape}の右で、直線がまず持ち上がったのと同じです）。
1000 ステップ後は $(1.0000, 1.0000)$ で、バイアスは $0.0000$、損失は $2\times10^{-14}$ です。
モデルは $1\times x_1 + 1\times x_2 + 0$、つまり\textbf{足し算そのもの}になりました。
学習データにない「$123+456$」を $(0.123, 0.456)$ として入れ、出力を 1000 倍すると、$579.0$ です（出力の最後の行）。
私たちは足し算のルールを書いていません。
用意したのは、例と、調整できる関数と、損失と、勾配降下法だけです。

ただし、この実験を言語モデルと同じだと思ってはいけません。
ここでは、数を数として入力し、足し算を表せる形のモデルを人が選びました。
言語モデルは、「\code{12+35=}」という\textbf{文字の並び}から、同じことを学ばなければなりません。
そのためには、はるかに柔軟なモデルが必要です。
それを作るのが、次章から第 III 部までの仕事です。

最後に、データを 100 や 1000 で割った理由を説明します。
勾配は「誤差 $\times$ 入力」でした。
割らずに使うと、入力も誤差も 1000 倍になり、勾配は 100 万倍になります。
すると、同じ学習率では 1 回の更新の幅が大きすぎて発散します（演習 \ref{ex:ml-scale}）。
入力と正解を 1 前後の大きさにそろえることは、学習を安定させる基本の工夫です。

この節では、勾配の計算を自動微分に任せました。
人が考えるのは「モデルと損失をどう計算するか」だけになり、同じ 4 手順のループで足し算さえ例から学べました。

\section{なぜ二乗誤差なのか：最尤推定という見方}\label{sec:ml-mle}
\skimread{初読では、結論だけ押さえれば先へ進めます。結論は「二乗誤差は、ずれが正規分布に従うと仮定した最尤推定から出てくる」（定理\ref{thm:ml-mse-mle}）です。
節末の囲み「損失関数の作り方」のレシピは次節で使うので、囲みと「ここまでの道のり」は読んでください。}

\secref{sec:ml-supervised}では、「打ち消し合わないから」という理由で、誤差を 2 乗しました。
しかし、それなら 4 乗でも絶対値でもよいはずです。
この節の問いは、「損失関数を決める原理はあるのか」です。
答えは、\chapref{ch:math}で学んだ\textbf{最尤推定}です。
この原理は、次節で分類の損失を作るときにも、\chapref{ch:lm}で言語モデルの損失を作るときにも使います。

最尤推定は、「観測されたデータが最も起こりやすくなるように、パラメータを選ぶ」という考え方でした。
コインを 10 回投げて表が 7 回出たなら、表の確率は 0.7 と考えるのが最も自然でした。
データが起こる確率を、パラメータの関数として見たものが\textbf{尤度}（ゆうど）でした。

導出に入る前に、この節の議論の骨組みを示します（\figref{fig:ml-mle-skeleton}）。
(a) 予測を、「一つの値」ではなく、正解 $y$ の確率分布 $p(y\mid x)$ として作ります（最尤推定は、モデルが確率を出さなければ使えません）。
(b) その分布の形に、直線の値 $wx+b$ を中心とする正規分布を選びます。
(c) この分布のもとで最尤推定をすると、(d) 二乗誤差を最小にする $w$ と $b$ が得られます。
(a) と (b) は、データから導かれることではなく、モデルを作る人が選ぶ\textbf{設計}（仮定）です。
(c) は計算の手順で、(d) だけが、そこから導かれる\textbf{結論}です。

\begin{figure}[tbp]
\centering
\begin{tikzpicture}[node distance=5mm,
  st/.style={minimum width=27mm,minimum height=14mm,text width=25mm,inner sep=1mm,font=\footnotesize}]
  \node[orangebox,st] (a) {\textbf{(a) 設計}\\予測を、値の確率分布\\$p(y\mid x)$ として作る};
  \node[orangebox,st,right=of a] (b) {\textbf{(b) 設計}\\分布の形に\\正規分布を選ぶ};
  \node[box,st,right=of b] (c) {\textbf{(c) 手順}\\最尤推定をする\\（尤度を最大にする）};
  \node[greenbox,st,right=of c] (d) {\textbf{(d) 結論}\\二乗誤差が最小の\\$w,b$ と同じになる};
  \draw[arr] (a) -- (b); \draw[arr] (b) -- (c); \draw[arr] (c) -- (d);
  \node[lbl,anchor=north] at ($(a.south west)!0.5!(b.south east)+(0,-1mm)$) {\eqref{eq:ml-gauss-noise}、\eqref{eq:ml-gauss-density}};
  \node[lbl,anchor=north] at ($(c.south)+(0,-1mm)$) {\eqref{eq:ml-likelihood}};
  \node[lbl,anchor=north] at ($(d.south)+(0,-1mm)$) {\eqref{eq:ml-gauss-nll}、定理\ref{thm:ml-mse-mle}};
\end{tikzpicture}
\caption{この節の議論の骨組み。橙の 2 つは、人が選ぶ設計（仮定）である。
そのもとで最尤推定（青）をすると、二乗誤差の最小化（緑）が結論として導かれる。箱の下は、対応する式の番号。}
\label{fig:ml-mle-skeleton}
\end{figure}

まず、(a) と (b) を式にします。
直線 $wx+b$ は、数を一つ出すだけです。
一方、実際のデータは直線の上下にばらついていました（\figref{fig:ml-scatter}）。
同じ文字数の断片でも、トークン数は一つに決まりません。
そこで、このばらつきも含めてモデルにします。
「正解 $y$ は、直線の値に、偶然のずれ $\varepsilon$ が加わってできた」と仮定するのです。
\begin{equation}
y = wx + b + \varepsilon,\qquad \varepsilon \sim \calN(0, \sigma^2)
\label{eq:ml-gauss-noise}
\end{equation}
$\calN(0,\sigma^2)$ は、平均 0、標準偏差 $\sigma$ の正規分布です（\chapref{ch:math}）。
$\varepsilon \sim$ は「$\varepsilon$ がこの分布に従う」と読みます。
ずれは 0 の近くが最も起こりやすく、大きなずれほど起こりにくい、という仮定です。
$\varepsilon$ は「正解 $-$ 予測」で、\secref{sec:ml-supervised}の誤差とは符号が逆です（2 乗するので、結論は同じです）。
$\sigma$ がどの入力 $x$ でも同じで、ずれがデータごとに無関係（独立）であることも仮定します。
$\sigma$ は、学習で求めず、人が置く数として扱います。

この仮定のもとでは、入力 $x$ を決めても $y$ は一つに決まりません。
$y$ は、直線の値 $wx+b$ を中心とする正規分布に従います（\figref{fig:ml-gauss-noise}の左）。
モデルの出力は、一つの値 $wx+b$ から、$y$ の確率分布 $p(y\mid x)$ に変わりました。
その確率密度は、\chapref{ch:math}の正規分布の式で、平均を $wx+b$ に置き換えたものです。
\begin{equation}
p(y \mid x) = \frac{1}{\sqrt{2\pi}\,\sigma}\exp\!\left(-\frac{\bigl(y - (wx+b)\bigr)^2}{2\sigma^2}\right)
\label{eq:ml-gauss-density}
\end{equation}
この式は、何かから導いたものではありません。
「入力が $x$ なら、正解 $y$ はこのあたりになりやすい」という予測の分布を、(a) と (b) の設計どおりに書いたものです。
\chapref{ch:math}では、係数を $\frac{1}{\sqrt{2\pi\sigma^2}}$ と書きました。
$\sqrt{2\pi\sigma^2}=\sqrt{2\pi}\,\sigma$ なので、同じ式です。
正解 $y$ が分布の中心 $wx+b$ に近いほど、この値は大きくなります。
指数関数の中身が 0 に近づき、$\exp$ の値が最大の 1 に近づくからです。

\begin{figure}[tbp]
\centering
\begin{tikzpicture}
\begin{axis}[width=0.5\linewidth,height=0.315\linewidth,xlabel={文字数},ylabel={トークン数},
  xmin=0,xmax=210,ymin=-8,ymax=135,ytick={0,50,100},
  tick label style={font=\scriptsize},label style={font=\footnotesize}]
\addplot[bookaccent,thick,domain=0:210,samples=2] {0.6*x};
\foreach \c/\m in {40/24,80/48,160/96} {
  \edef\temp{\noexpand\addplot[bookgreen,thick,fill=figgreen,fill opacity=0.6,domain=\m-30:\m+30,samples=41,variable=\noexpand\t]
    ({\c+28*exp(-(t-\m)^2/200)},{t}) -- cycle;
    \noexpand\draw[black!40,thin] (axis cs:\c,\m-30) -- (axis cs:\c,\m+30);}
  \temp
}
\addplot[only marks,mark=*,mark size=1.8pt,bookblue] coordinates {(40,21) (80,42) (160,101)};
\node[font=\scriptsize,text=bookaccent,anchor=south east] at (axis cs:136,92) {$w=0.6$};
\end{axis}
\end{tikzpicture}\hfill
\begin{tikzpicture}
\begin{axis}[width=0.5\linewidth,height=0.315\linewidth,xlabel={ずれ（正解 $-$ 予測、トークン）},ylabel={確率密度},
  xmin=-42,xmax=42,ymin=0,ymax=0.06,scaled y ticks=false,ytick={0,0.02,0.04,0.06},
  tick label style={font=\scriptsize,/pgf/number format/fixed,/pgf/number format/precision=2},
  label style={font=\footnotesize}]
\addplot[ybar,bar width=4,fill=figblue,draw=bookblue] table[col sep=comma,x=center,y=density]{code/ch03/out/resid_hist.csv};
\addplot[bookaccent,thick,domain=-42:42,samples=80] {0.04182*exp(-x*x/182.02)};
\end{axis}
\end{tikzpicture}
\caption{左：ばらつきを含めたモデル。入力 $x$ ごとに、直線の値（$w=0.6$、$b=0$）を中心とする正規分布を考える。
緑の山は、3 つの文字数（40、80、160 字）での分布（$\sigma=10$ トークン）を横向きに描いたもので、右への出っ張りが大きいトークン数ほど起こりやすい（出っ張りの幅は、横軸の文字数とは無関係）。
青い点が実際のデータ（3 件）で、山の中心（直線の値）に近いほど「起こりやすい」。
右：断片を 1 万件に増やして測った、直線からのずれ（正解 $-$ 予測）の分布（青い棒）と、同じ標準偏差の正規分布（赤い線）。}
\label{fig:ml-gauss-noise}
\end{figure}

次は (c) です。
データは $N$ 件あります。
それぞれのずれは独立だと仮定したので、全部が同時に起こる確率は掛け算になります。
$y$ は連続的な値として扱うので、確率の代わりに確率密度（分布の山の高さ）を掛け、これを尤度と呼びます。
知りたいのは、どのパラメータで尤度が大きいかという大小だけなので、密度で足ります。
\begin{equation}
\text{尤度} = \prod_{i=1}^{N} p\bigl(y^{(i)} \mid x^{(i)}\bigr)
\label{eq:ml-likelihood}
\end{equation}
最尤推定は、この尤度が最大になる $w$ と $b$ を選びます。
\secref{sec:ml-supervised}の 3 件のデータで、実際に計算します。
ここだけは、単位を百トークンから元のトークンに戻し、$b=0$ とします。
ばらつきの幅 $\sigma$ は、仮に 10 トークンと置きます（\secref{sec:ml-gd}で見た典型的な誤差の大きさです）。
すぐあとで分かるとおり、$\sigma$ をいくつに置いても、最良の $w$ は変わりません。
3 通りの $w$ について、3 件それぞれの確率密度と、その積（尤度）を計算しました（\file{code/ch03/gauss_mle.py}）。
尤度の対数を取って符号を変えたものも計算します。
これを\term{負の対数尤度}{negative log-likelihood, NLL}と呼びます。
右端の列は、誤差の 2 乗の和です（\tabref{tab:ml-hand-loss}の「誤差$^2$」の 3 件の和を、トークン単位に戻すため $100^2$ 倍した値）。

\outputfile{ch03/out/mle_table.txt}

$w=0.6$ の行を読みます。
1 件目（40 字）の確率密度を、\eqref{eq:ml-gauss-density}で手で計算してみます。
予測は $0.6\times40=24$ トークン、正解は 21 トークンなので、ずれは $21-24=-3$ です。
係数は $\frac{1}{\sqrt{2\pi}\times10}=0.0399$ です。
指数関数の部分は $\exp\bigl(-\frac{(-3)^2}{2\times10^2}\bigr)=\exp(-0.045)=0.956$ です。
掛けると $0.0399\times0.956=0.0381$ で、出力の値と一致します。
2 件目（ずれ $-6$）と 3 件目（ずれ $+5$）も同じ計算で、$0.0333$ と $0.0352$ になります。
3 つを掛けると $4.47\times10^{-5}$ で、これが尤度です。
右端の列の二乗誤差の和は、$3^2+6^2+5^2=70$ です。
この尤度は、3 つの $w$ の中で最大です。
負の対数尤度は最小で、二乗誤差の和も最小です。
最後の 2 行は、$w$ を $0.001$ 刻みで調べた結果です。
尤度が最大になる $w$ と、二乗誤差が最小になる $w$ は、どちらも $0.606$ でした。
これは偶然ではありません。

\subsection*{対数を取ると、二乗誤差が現れる}

最後に、(d) の結論を導きます。
尤度は掛け算なので、件数が増えると、1 より小さい数を何度も掛けることになります。
すると値が小さくなりすぎて、計算機の上では 0 になってしまいます（\chapref{ch:math}のアンダーフロー）。
そこで対数を取ります。
対数は掛け算を足し算に変え、しかも大小関係を変えません（\chapref{ch:math}）。
尤度が最大になる場所と、対数尤度が最大になる場所は同じです。

1 件分の対数を計算します。
\eqref{eq:ml-gauss-density}は「係数 $\times$ 指数関数」の形です。
積の対数は対数の和で、$\log\exp(a) = a$ なので、次のようになります。
\begin{equation}
\log p(y \mid x) = -\log\bigl(\sqrt{2\pi}\,\sigma\bigr) - \frac{\bigl(y - (wx+b)\bigr)^2}{2\sigma^2}
\end{equation}
\textbf{2 乗が現れました。}正規分布の式の、指数関数の中にあった 2 乗です。
これを $N$ 件分足し、符号を変えると、負の対数尤度になります。
\begin{equation}
-\sum_{i=1}^{N}\log p\bigl(y^{(i)} \mid x^{(i)}\bigr)
= \underbrace{N\log\bigl(\sqrt{2\pi}\,\sigma\bigr)}_{w,\,b\ \text{に関係しない}}
+ \frac{1}{2\sigma^2}\underbrace{\sum_{i=1}^{N}\bigl(y^{(i)} - (wx^{(i)}+b)\bigr)^2}_{\text{二乗誤差の和}}
\label{eq:ml-gauss-nll}
\end{equation}
第 1 項は $w$ も $b$ も含まないので、どの直線を選んでも同じ値です。
第 2 項は、二乗誤差の和に正の定数 $\frac{1}{2\sigma^2}$ を掛けたものです。
したがって、$\sigma$ がいくつであっても、負の対数尤度を最小にする $(w,b)$ は、二乗誤差を最小にする $(w,b)$ と同じです。
これが (d) の結論です。
和を $N$ で割って平均にしても、最小になる場所は変わらないので、平均二乗誤差（\eqref{eq:ml-mse}）でも同じです。
数値でも確かめましょう。
$w=0.6$ では、第 1 項が $3\times\log(\sqrt{2\pi}\times 10) = 3\times3.2215 = 9.665$ です。
第 2 項は、$2\sigma^2 = 2\times10^2 = 200$ なので、$70/200 = 0.35$ です。
合計は $10.015$ で、出力の値と一致します。
% [figpass] 案 3・4
\figref{fig:ml-mle-numbers}の上は、ここまでの $w=0.6$ の計算（3 件の確率密度 → 尤度 → 負の対数尤度 → 二つの項への分解）を、左から順に並べたものです。
\figref{fig:ml-mle-numbers}の下は、$\sigma$ を 5、10、20 トークンの 3 通りに変えて、3 件のデータの尤度を $w$ の関数として描いたものです（山の高さの最大を 1 にそろえた）。
$\sigma$ を変えると山の幅は変わりますが、最大になる位置は 0.606 のまま動きません。
係数 $\frac{1}{2\sigma^2}$ は山の鋭さを決めるだけで、どの $w$ で最大になるかは、二乗誤差の和だけで決まるからです。

\begin{figure}[tbp]
\centering
\begin{tikzpicture}[x=1mm,y=1mm,
  c1/.style={plainbox,font=\scriptsize,minimum width=32mm,minimum height=5.5mm,inner sep=1pt},
  c2/.style={plainbox,font=\scriptsize,minimum width=13mm,minimum height=5.5mm,inner sep=1pt},
  c3/.style={greenbox,font=\scriptsize,minimum width=17mm,minimum height=5.5mm,inner sep=1pt},
  big/.style={font=\footnotesize,minimum width=21mm,minimum height=11mm,align=center,inner sep=1pt},
  sm/.style={font=\scriptsize,text=bookgray,inner sep=1pt}]
  \foreach \yy/\a/\d/\pd [count=\k] in {0/{40 字：予測 24、正解 21}/{$-3$}/{0.0381},
      -7/{80 字：予測 48、正解 42}/{$-6$}/{0.0333},
      -14/{160 字：予測 96、正解 101}/{$+5$}/{0.0352}} {
    \node[c1] (a\k) at (16,\yy) {\a};
    \node[c2] (d\k) at (42,\yy) {ずれ \d};
    \node[c3] (p\k) at (61.5,\yy) {密度 \pd};
    \draw[arr] (a\k) -- (d\k); \draw[arr] (d\k) -- (p\k);
  }
  \node[sm,anchor=south] at (42,3) {正解 $-$ 予測};
  \node[sm,anchor=south] at (61.5,3) {\eqref{eq:ml-gauss-density}、$\sigma=10$};
  \draw[decorate,decoration={brace,amplitude=4pt}] (71,2.8) -- (71,-16.8);
  \node[box,big] (lik) at (86,-7) {尤度（3 つの積）\\$4.47\times10^{-5}$};
  \draw[arr] (73.5,-7) -- (lik.west);
  \node[redbox,big] (nll) at (117,-7) {負の対数尤度\\$10.015$};
  \draw[arr] (lik) -- node[above,sm]{$-\log$} (nll);
  \node[font=\footnotesize,anchor=east] at (97,-16) {$10.015\ =$};
  \node[font=\footnotesize] (t1) at (102,-16) {$9.665$};
  \node[font=\footnotesize] at (108,-16) {$+$};
  \node[font=\footnotesize] (t2) at (118,-16) {$0.35$};
  \node[sm,anchor=north,align=center,xshift=-2.5mm] at (t1.south) {$w,b$ に\\関係しない};
  \node[sm,anchor=north,align=center] at (t2.south) {二乗誤差の和 $\div\,2\sigma^2$\\$=70\div(2\times10^2)$};
\end{tikzpicture}

\vspace{2mm}
\begin{tikzpicture}
\begin{axis}[width=0.66\linewidth,height=0.3\linewidth,xlabel={重み $w$},ylabel={尤度（最大を 1 にそろえた）},
  xmin=0.35,xmax=0.85,ymin=0,ymax=1.22,ytick={0,0.5,1},
  tick label style={font=\scriptsize},label style={font=\footnotesize},
  legend style={font=\scriptsize,at={(1.03,1)},anchor=north west,draw=black!30},legend cell align=left]
% 3 件のデータの二乗誤差の和は 33600 w^2 - 40720 w + 12406（トークン単位）
\addplot[bookaccent,very thick,domain=0.35:0.85,samples=121] {exp(-33600*(x-0.605952)^2/(2*5^2))};
\addplot[bookblue,thick,densely dashed,domain=0.35:0.85,samples=121] {exp(-33600*(x-0.605952)^2/(2*10^2))};
\addplot[bookgreen,thick,dashdotted,domain=0.35:0.85,samples=121] {exp(-33600*(x-0.605952)^2/(2*20^2))};
\legend{$\sigma=5$,$\sigma=10$,$\sigma=20$}
\draw[black!50,thin] (axis cs:0.606,0) -- (axis cs:0.606,1.22);
\node[font=\scriptsize,anchor=south east,inner sep=1pt] at (axis cs:0.603,1.03) {$w=0.606$};
\end{axis}
\end{tikzpicture}
\caption{最尤推定の計算（数値例）。上：$w=0.6$、$\sigma=10$ トークンのときの計算の順序（出力 \file{mle_table.txt} の $w=0.6$ の行）。
3 件のずれから確率密度を求め、掛けて尤度に、負の対数を取って負の対数尤度にする。負の対数尤度は、$w,b$ に関係しない項と、二乗誤差の和に比例する項に分かれる。
下：$\sigma$ を 3 通りに変えたときの、3 件のデータの尤度（\eqref{eq:ml-gauss-nll}から計算）。山の幅は $\sigma$ で変わるが、最大になる位置はどれも $w=0.606$ にある。}
\label{fig:ml-mle-numbers}
\end{figure}
% [/figpass]

\begin{theorem}[二乗誤差と最尤推定]\label{thm:ml-mse-mle}
正解が「モデルの予測 $+$ 正規分布に従うずれ」でできていると仮定する。
ずれの標準偏差 $\sigma$ はどの入力でも同じで、ずれはデータごとに独立であるとする。
このとき、尤度を最大にすること（最尤推定）と、二乗誤差を最小にすることは、同じパラメータを与える。
\end{theorem}

二乗誤差は、気まぐれに選んだものではありませんでした。
「ずれは正規分布に従う」という仮定を置いて最尤推定をすると、計算の結果として出てくる損失だったのです。

では、仮定は、このデータでどこまで成り立つのでしょうか。
\figref{fig:ml-gauss-noise}の右は、同じ作り方で断片を 1 万件に増やし、直線からのずれを実際に調べたものです。
山の形は正規分布に似ていますが、中央がやや尖っています。
「どの入力でも同じ $\sigma$」も正確ではなく、文字数が多いほど、点の上下の広がりは大きくなります（\figref{fig:ml-scatter}の左）。
仮定は、おおよそ成り立つ近似です。
仮定が正確でなくても、二乗誤差を最小にする計算はでき、得られる直線も変わりません。
変わるのは、「その直線が最尤推定と一致する」という意味づけです。
ずれの分布が正規分布から大きく外れる問題では、別の損失のほうが合います（演習 \ref{ex:ml-laplace}）。

\begin{intuition}{損失関数の作り方}
この節の結論は、一つのレシピにまとめられます。
(1) モデルを、正解についての\textbf{確率分布}を出力するように作る。分布の形は、人が選ぶ。
(2) その分布のもとで、実際の正解が起こる確率の対数を取り、符号を変えたもの（負の対数尤度）を損失にする。
正解が数なら、(1) で正規分布を選ぶと、損失は二乗誤差になります。
正解が「種類」のときに同じレシピを使うと何が出てくるかを、次の節で見ます。
\end{intuition}

\subsection*{ここまでの道のり}

章の半分を過ぎました。ここまでを 3 行でまとめます。
\begin{itemize}
\item 機械学習は、データ・モデル・損失関数・最適化の 4 要素でできています。
\item 線形回帰では、モデルは直線、損失は二乗誤差で、勾配降下法によって損失を小さくしました。
\item 損失は「負の対数尤度」として作れます。二乗誤差は、その一例でした。
\end{itemize}
残る話題は三つです。分類、データが大きいときの工夫、そして「学習に使っていないデータでも当たるか」です。

\section{分類：確率を出力するモデル}\label{sec:ml-classification}

ここまでの問題では、正解は数でした。
このような問題を\term{回帰}{regression}と呼びます。
一方、言語モデルが当てるのは、4096 種類のトークンのうちの\textbf{どれか}です。
正解が「種類」である問題を\term{分類}{classification}と呼び、種類のことを\term{クラス}{class}と呼びます。
この節の問いは、「正解が種類のとき、モデルと損失をどう作るか」です。
4 要素のうち、取り替えるのはモデルと損失関数だけで、作り方には前節のレシピがそのまま使えます。
レシピは、(1) モデルに、正解についての確率分布を出力させ、(2) 実際の正解が起こる確率の負の対数（負の対数尤度）を損失にする、の二つでした。

\subsection*{2 クラスの分類：ロジスティック回帰}

いちばん簡単な、2 択の問題から始めます。
同じ断片のデータを使い、問いを変えます。
「この断片は 60 トークンに収まるか」を、文字数から当てるのです。
収まるなら $y=1$、収まらないなら $y=0$ とします。
100 件のうち、$y=1$ は 42 件でした。

レシピの (1) に従い、モデルには「$y=1$ である確率 $p$」を出力させます。
直線の値 $z = wx+b$ は、負にも、1 より大きくもなるので、そのままでは確率になりません。
そこで、どんな実数も 0 と 1 の間の値に変える関数を通します。
\begin{equation}
p = \sigmoid(z) = \frac{1}{1+e^{-z}},\qquad z = wx+b
\label{eq:ml-sigmoid}
\end{equation}
この $\sigmoid$ を\term{シグモイド関数}{sigmoid function}と呼びます（前節の標準偏差 $\sigma$ とは別物）。
\figref{fig:ml-sigmoid}の左がグラフです。
$z=0$ で $0.5$、$z$ が大きいと 1 に、小さいと 0 に近づきます。
確率に直す前の値 $z$ を\term{ロジット}{logits}と呼びます。
このモデルを\term{ロジスティック回帰}{logistic regression}と呼びます。
名前に「回帰」とありますが、分類のモデルです。

\begin{figure}[tbp]
\centering
\pgfplotsset{mlthree/.style={width=0.385\linewidth,height=0.29\linewidth,
  title style={font=\footnotesize,yshift=-1.5mm},tick label style={font=\scriptsize},
  label style={font=\scriptsize}}}
\begin{tikzpicture}
\begin{axis}[mlthree,title={シグモイド $\sigmoid(z)$},xlabel={ロジット $z$},
  xmin=-6,xmax=6,ymin=-0.05,ymax=1.05,ytick={0,0.5,1}]
\addplot[bookblue,very thick,domain=-6:6,samples=80] {1/(1+exp(-x))};
\addplot[only marks,mark=*,mark size=1.5pt,bookaccent] coordinates {(0,0.5) (2,0.881) (-2,0.119)};
\node[font=\scriptsize,anchor=north west] at (axis cs:2,0.86) {$0.88$};
\node[font=\scriptsize,anchor=south east] at (axis cs:-2,0.14) {$0.12$};
\end{axis}
\end{tikzpicture}\hfill
\begin{tikzpicture}
\begin{axis}[mlthree,title={損失（負の対数）},xlabel={正解に割り当てた確率},
  xmin=0,xmax=1,ymin=0,ymax=5]
\addplot[bookaccent,very thick,domain=0.006:1,samples=120] {-ln(x)};
\addplot[only marks,mark=*,mark size=1.5pt,bookblue] coordinates {(0.9,0.105) (0.5,0.693) (0.1,2.303)};
\node[font=\scriptsize,anchor=south west] at (axis cs:0.1,2.303) {$2.30$};
\node[font=\scriptsize,anchor=south west] at (axis cs:0.5,0.693) {$0.69$};
\node[font=\scriptsize,anchor=south] at (axis cs:0.88,0.15) {$0.11$};
\end{axis}
\end{tikzpicture}\hfill
\begin{tikzpicture}
\begin{axis}[mlthree,title={学習後の確率 $p$},xlabel={文字数},
  xmin=10,xmax=210,ymin=-0.05,ymax=1.05,ytick={0,0.5,1}]
\draw[black!50,densely dashed] (axis cs:98.2,-0.05) -- (axis cs:98.2,1.05);
\addplot[only marks,mark=*,mark size=0.9pt,bookblue!70] table[col sep=comma,x=chars,y=y]
  {code/ch03/out/logistic_points.csv};
\addplot[bookaccent,very thick] table[col sep=comma,x=chars,y=p]{code/ch03/out/logistic_curve.csv};
\node[font=\scriptsize,anchor=west] at (axis cs:101,0.62) {98.2 字};
\end{axis}
\end{tikzpicture}
\caption{縦軸は、どの図も、上の題に書いた量。
左：シグモイド関数（数値は $z=\pm 2$ での値）。
中央：1 件の損失。横軸は、モデルが正解に割り当てた確率（$y=1$ なら $p$、$y=0$ なら $1-p$）で、数値は確率 $0.1$、$0.5$、$0.9$ での損失。確率が低いと、損失は急に大きくなる。
右：学習後のロジスティック回帰。青い点は訓練データ（60 トークンに収まれば 1、収まらなければ 0）、赤い曲線はモデルが出す「収まる確率」$p$。
破線は、$p$ が $0.5$ をまたぐ文字数（98.2 字）。}
\label{fig:ml-sigmoid}
\end{figure}

次にレシピの (2) です。
正解が起こる確率は、$y=1$ なら $p$、$y=0$ なら $1-p$ です。
その対数の符号を変えたものが損失です。
二つの場合は、次の式にまとめられます。
\begin{equation}
\ell = -\bigl\{\,y\log p + (1-y)\log(1-p)\,\bigr\}
\label{eq:ml-bce}
\end{equation}
$y=1$ なら第 2 項が消えて $-\log p$ が残り、$y=0$ なら第 1 項が消えて $-\log(1-p)$ が残ります。
どちらの場合も、損失は「モデルが正解に割り当てた確率」の負の対数です。
この損失は、あとで定義する交差エントロピー損失（定義\ref{def:ml-cross-entropy}）の、クラスが二つの場合にあたります。
\figref{fig:ml-sigmoid}の中央は、正解に割り当てた確率を横軸に、損失を縦軸に取ったグラフです。
正解に確率 $0.9$ を割り当てていれば、損失は $0.11$ です。
$0.1$ しか割り当てていなければ $2.30$ で、0 に近づくと限りなく大きくなります。
正解に低い確率しか割り当てないほど、つまり自信を持って間違えるほど、損失は大きくなります。

あとは勾配降下法です。そのために、損失を $w$ と $b$ で微分します。
損失は、$w, b \to z \to p \to \ell$ という 3 段の合成なので、後ろ（$\ell$ の側）から 1 段ずつ微分して掛け合わせます（合成関数の微分）。
まず、\eqref{eq:ml-bce}を $p$ で微分します。
$\log p$ の微分は $\frac{1}{p}$ です。
$\log(1-p)$ の微分は、中身 $1-p$ の微分 $-1$ が掛かって、$-\frac{1}{1-p}$ です。
この二つを使い、通分すると次の形になります。
\begin{equation}
\pd{\ell}{p} = -\frac{y}{p} + \frac{1-y}{1-p} = \frac{-y(1-p) + (1-y)p}{p(1-p)} = \frac{p-y}{p(1-p)}
\end{equation}
次に、シグモイド関数 $p=\sigmoid(z)$ を $z$ で微分します。結果は $\dd{p}{z} = p(1-p)$ という簡単な形です。
筋道だけ示します。
$u=1+e^{-z}$ と置くと $p=\frac{1}{u}$ で、$\frac{1}{u}$ の微分は $-\frac{1}{u^2}$、$u$ の微分は $-e^{-z}$ です。
掛け合わせると、$\dd{p}{z}=\frac{e^{-z}}{(1+e^{-z})^2}=\frac{1}{1+e^{-z}}\times\frac{e^{-z}}{1+e^{-z}}$ です。
前の分数は $p$、後ろの分数は $1-p$ なので、$p(1-p)$ になります（演習 \ref{ex:ml-sigmoid}で、自分の手でも確かめてください）。
たとえば $z=0$ では $p=0.5$ なので、傾きは $0.5\times0.5=0.25$ です。
\figref{fig:ml-sigmoid}の左の曲線が、いちばん急な場所です。

$\pd{\ell}{p}$ と $\dd{p}{z}$ の二つを掛けると、分母がちょうど打ち消されます。
\begin{equation}
\pd{\ell}{z} = \pd{\ell}{p}\,\dd{p}{z} = \frac{p-y}{p(1-p)}\times p(1-p) = p - y
\label{eq:ml-logistic-dz}
\end{equation}
最後に、$z = wx+b$ を $w$ で微分すると $x$、$b$ で微分すると 1 です。これらを掛けて、勾配が得られます。
\begin{equation}
\pd{\ell}{w} = (p - y)\,x,\qquad \pd{\ell}{b} = p - y
\label{eq:ml-logistic-grad}
\end{equation}
線形回帰と同じ、「誤差 $\times$ 入力」の形です。ここでの誤差は「予測した確率 $-$ 正解」です。

% [figpass] 案 1
\figref{fig:ml-logistic-graph}は、この 3 段の計算を、数値例（50 字の断片 $x=0.5$、$w=-2$、$b=2$、正解 $y=1$）で描いた図です。
箱の中の黒い数が、左から右へ順に計算した値です。
矢印のそばの枠のない小さな赤い数は、各段の微分です。
赤い枠で囲んだ数は、$\ell$ の側から 1 段ずつ、各段の微分を掛けていった値です。
$\pd{\ell}{p}=-1.368$ に $\dd{p}{z}=0.197$ を掛けた $-0.269$ が、$p-y=0.731-1$ にちょうど一致しています。

\begin{figure}[tbp]
\centering
\begin{tikzpicture}[x=1cm,y=1cm,
  inp/.style={plainbox,font=\footnotesize,minimum width=18mm,minimum height=7mm},
  stg/.style={box,font=\footnotesize,minimum width=20mm,minimum height=13mm,align=center},
  gr/.style={font=\footnotesize,text=bookaccent,align=center,inner sep=2pt,draw=bookaccent,rounded corners=1pt,fill=white},
  loc/.style={font=\scriptsize,text=bookaccent,align=center,inner sep=1pt}]
  \node[inp] (w) at (0,1.2) {重み $w=-2$};
  \node[inp] (x) at (0,0) {入力 $x=0.5$};
  \node[inp] (b) at (0,-1.2) {バイアス $b=2$};
  \node[stg] (z) at (3.4,0) {ロジット\\$z=wx+b$\\$=1.00$};
  \node[stg] (p) at (6.9,0) {確率\\$p=\sigmoid(z)$\\$=0.731$};
  \node[redbox,stg] (l) at (10.4,0) {損失\\$\ell=-\log p$\\$=0.313$};
  \node[orangebox,font=\footnotesize,minimum height=6mm] (y) at (10.4,1.6) {正解 $y=1$};
  \draw[arr] (w.east) -- (z.west); \draw[arr] (x) -- (z); \draw[arr] (b.east) -- (z.west);
  \draw[arr] (z) -- (p); \draw[arr] (p) -- (l); \draw[arr] (y) -- (l);
  % 各段の微分（箱と箱の間の上）
  \node[loc,anchor=south] at (5.15,0.72) {$\dd{p}{z}=p(1-p)$\\$=0.197$};
  \node[loc,anchor=south] at (8.65,0.72) {$\pd{\ell}{p}=-\frac{1}{p}$\\[1.2mm]$=-1.368$};
  \node[loc,anchor=south west] at (1.3,1.02) {$\pd{z}{w}=x$};
  \node[loc,anchor=north west] at (1.3,-1.02) {$\pd{z}{b}=1$};
  % 損失の側から掛け合わせた値（箱の下・入力の上下）
  \node[gr,anchor=north] at (6.9,-0.8) {$\pd{\ell}{p}=-1.368$};
  \node[gr,anchor=north west] at (2.5,-0.8) {$\pd{\ell}{z}=-1.368\times0.197$\\$=-0.269\ (=p-y)$};
  \node[gr,anchor=south west] at (-0.9,1.7) {$\pd{\ell}{w}=-0.269\times0.5=-0.134$};
  \node[gr,anchor=north west] at (-0.9,-1.7) {$\pd{\ell}{b}=-0.269\times1=-0.269$};
  % 逆向きの矢印
  \draw[->,thick,dashed,bookaccent] (10.4,-2.3) -- node[below,loc]{損失の側から、各段の微分を掛けていく（合成関数の微分）} (3.4,-2.3);
\end{tikzpicture}
\caption{ロジスティック回帰の損失を、計算の段ごとに描いた図（計算グラフ。数値例）。黒い矢印と箱の中の数は、入力から損失への計算。
枠のない小さな赤字は各段の微分（箱と箱をつなぐ矢印のそば）、赤い枠で囲んだ値は、それを損失の側から掛け合わせた値。
$\pd{\ell}{p}=\frac{p-y}{p(1-p)}$（$y=1$ では $-\frac{1}{p}$）と $\dd{p}{z}=p(1-p)$ の積で分母が消え、ロジットでの微分は $p-y$ になる。
パラメータの勾配は、それに入力 $x$（$w$ の場合）か 1（$b$ の場合）を掛けたもの。}
\label{fig:ml-logistic-graph}
\end{figure}
% [/figpass]

これを実装して（\file{code/ch03/logistic_regression.py}）、学習率 $1.0$ で 5000 ステップ学習しました。
最初の損失は $0.693$ です。
$w=b=0$ では $z=0$ なので、モデルはどの断片にも $p=\sigmoid(0)=0.5$ を出し、損失が $-\log 0.5=0.693$ になるからです。
学習後の損失は、$0.077$ まで下がりました。
\figref{fig:ml-sigmoid}の右が、学習後のモデルです。
短い断片では確率がほぼ 1、長い断片ではほぼ 0 で、その間を S 字の曲線がつないでいます。
$0.5$ をまたぐ場所（破線）は 98.2 字で、\secref{sec:ml-gd}の「100 字で約 60 トークン」という結果と合っています。
確率が $0.5$ より大きければ「収まる」と予測することにすると、学習に使っていない別の 100 件の 96\% で当たりました。

\subsection*{多クラスの分類：softmax 回帰}

クラスが 3 つ以上あるときに進みます。クラスの数を $K$ と書きます。
進め方は、2 クラスのときと同じ 3 手順です。
確率を出力するモデルを作り、負の対数尤度を損失にし、その勾配を求めます。
例には、境界を目で見られるように、人工的に作った平面上の点 $\vx=(x_1,x_2)^\T$ を使います。
3 つの中心のまわりに 100 点ずつ散らした 300 件で、各点がどの中心のものかを当てます（点の散らばりは、あとの\figref{fig:ml-blobs}の中央）。

まず、クラスごとに「そのクラスらしさ」を表す数を、入力の重み付きの和で作ります。
この数を、\chapref{ch:math}と同じく\textbf{スコア}と呼びます。
クラス $k$ のスコアは $z_k = \vw_k^\T\vx + b_k$ です。
$\vw_k$ は\textbf{クラス $k$ の重みベクトル}、$b_k$ はクラス $k$ のバイアスで、\eqref{eq:ml-linear-multi}と同じ形の式を、クラスごとに一つずつ持ちます。
$K$ 個をまとめると、行列とベクトルの積で書けます。
\begin{equation}
\vz = \mW\vx + \vb,\qquad \mW\in\R^{K\times 2},\quad \vb\in\R^{K}
\label{eq:ml-logits}
\end{equation}
$\mW$ を\textbf{重み行列}と呼びます。
$k$ 行目に、クラス $k$ の重みベクトル $\vw_k$ の成分が横に並びます（\figref{fig:ml-weight-rows}）。
$\vb$ は、バイアス $b_k$ を縦に並べたベクトルです。
行列とベクトルの積の $k$ 番目の成分は、行列の $k$ 行目とベクトルの内積なので（\chapref{ch:math}）、$\mW\vx+\vb$ の $k$ 番目の成分が、クラス $k$ のスコア $z_k$ になります。
$K$ 個のスコアを縦に並べた $\vz$ が、softmax に入れる前の値、つまりロジットです。

\begin{figure}[tbp]
\centering
\begin{tikzpicture}[x=13mm,y=5.8mm,
  c/.style={cell,minimum width=13mm,minimum height=5.8mm},
  hl/.style={c,fill=figorange,draw=bookorange,thick},
  hd/.style={font=\footnotesize,anchor=south}]
  % 重み行列 W（3 行 2 列）
  \node[c] at (0,-1) {$-1.60$}; \node[c] at (1,-1) {$-0.85$};
  \node[c] at (0,-3) {$-0.16$}; \node[c] at (1,-3) {$1.86$};
  \node[hl] at (0,-2) {$1.76$}; \node[hl] at (1,-2) {$-1.01$};
  \node[hd] at (0.5,-0.45) {重み行列 $\mW$（3 行 2 列）};
  \foreach \r in {1,2,3} \node[lbl,anchor=east] at (-0.6,-\r) {$\r$ 行目（クラス $\r$）};
  % 入力 x
  \node[c] at (2.45,-1.5) {$2$}; \node[c] at (2.45,-2.5) {$-1$};
  \node[hd] at (2.45,-0.45) {入力 $\vx$};
  \node[font=\small] at (3.3,-2) {$+$};
  % バイアス b
  \node[c] at (4.15,-1) {$-0.05$}; \node[hl] at (4.15,-2) {$0.15$}; \node[c] at (4.15,-3) {$-0.10$};
  \node[hd] at (4.15,-0.45) {バイアス $\vb$};
  \node[font=\small] at (5.0,-2) {$=$};
  % スコア z
  \node[c] at (5.85,-1) {$-2.40$}; \node[hl] at (5.85,-2) {$4.68$}; \node[c] at (5.85,-3) {$-2.28$};
  \node[hd] at (5.85,-0.45) {スコア $\vz$};
\end{tikzpicture}
\caption{重み行列 $\mW$ の $k$ 行目が、クラス $k$ の重みベクトル $\vw_k$（を横に並べた $\vw_k^\T$）である。
数値は、この節の最後に学習して得られる値（小数第 2 位まで）で、入力は例として $\vx=(2,-1)^\T$ とした。
橙の 2 行目 $\vw_2^\T=(1.76,\ -1.01)$ と $\vx$ の内積にバイアス $b_2$ を足すと、クラス 2 のスコア $z_2=1.76\times 2+(-1.01)\times(-1)+0.15=4.68$ になる。ほかの行も同様。}
\label{fig:ml-weight-rows}
\end{figure}

次に、ロジットを確率に直します。
使うのは、\chapref{ch:math}で学んだ softmax 関数です。
\begin{equation}
p_k = \frac{e^{z_k}}{\sum_{j=1}^{K} e^{z_j}}
\label{eq:ml-softmax}
\end{equation}
各 $p_k$ は正で、合計は 1 になります。
たとえば $\vz=(2,1,0)$ で計算してみます（\chapref{ch:math}と同じ例。成分を横に並べて書きますが、列ベクトルです）。
$e^2=7.389$、$e^1=2.718$、$e^0=1$ で、合計は $11.107$ です。
それぞれを合計で割ると、$\vp=(7.389,\ 2.718,\ 1)\div 11.107=(0.665, 0.245, 0.090)$ です。
スコアが 1 だけ高いと、確率は $e\approx2.7$ 倍です。
このモデルを \textbf{softmax 回帰}（softmax regression）と呼びます。
$K=2$ のときは、ロジスティック回帰と同じものになります（演習 \ref{ex:ml-sigmoid}）。

正解の書き方も決めておきます。
正解がクラス 2 なら、$\vy=(0,1,0)$ と書きます。
正解のクラスの成分だけが 1 で、ほかは 0 のベクトルです。
この書き方を \textbf{one-hot 表現}と呼びます。
これは、「クラス 2 である確率が 1」という確率分布と見ることもできます。

損失は、レシピどおり、正解のクラスに割り当てた確率の負の対数です。
正解のクラスを $c$ とすると $\ell = -\log p_c$ です。
one-hot 表現を使うと、和の形でも書けます。
$y_k$ は $k=c$ のときだけ 1 なので、和を取っても $-\log p_c$ だけが残るからです。

\begin{definition}[交差エントロピー損失]\label{def:ml-cross-entropy}
正解のクラスを $c$、その one-hot 表現を $\vy$、モデルが出力した確率を $\vp$ とする。
\begin{equation}
\ell = -\log p_c = -\sum_{k=1}^{K} y_k \log p_k
\label{eq:ml-cross-entropy}
\end{equation}
を\textbf{交差エントロピー損失}（cross-entropy loss）と呼ぶ。
これは、正解のクラスについての負の対数尤度である。
\end{definition}

右辺は、\chapref{ch:math}で学んだ交差エントロピー $H(\vy,\vp)$ の式そのものです。
分類の問題で、正解を one-hot 表現で与えるとき、「交差エントロピー損失」と「負の対数尤度」は同じものになります。
負の対数尤度のほうが広い言葉で、前節の二乗誤差も、正規分布を選んだときの負の対数尤度でした。
先ほどの例で正解がクラス 2 なら、損失は $-\log 0.245 = 1.41$ です。
\figref{fig:ml-softmax-flow}に、入力から損失までの計算の順序をまとめました。

\begin{figure}[tbp]
\centering
\begin{tikzpicture}[node distance=5mm and 10mm,
  vec/.style={plainbox,font=\footnotesize,minimum height=14mm,minimum width=13mm}]
  \node[vec] (x) {入力 $\vx$};
  \node[box,right=of x,font=\footnotesize,minimum height=10mm] (lin) {$\mW\vx+\vb$\\重み行列 $\mW$\\バイアス $\vb$};
  \node[vec,right=of lin] (z) {ロジット $\vz$\\$2.0$\\$1.0$\\$0.0$};
  \node[box,right=of z,font=\footnotesize,minimum height=10mm] (sm) {softmax};
  \node[vec,right=of sm,fill=figgreen,draw=bookgreen] (p) {確率 $\vp$\\$0.665$\\$0.245$\\$0.090$};
  \node[vec,below=9mm of p,fill=figorange,draw=bookorange] (y) {正解 $\vy$\\$0$\\$1$\\$0$};
  % 損失の箱は softmax の真下、勾配の箱はロジットの真下に置く（破線の矢印をまっすぐにするため）
  \node[redbox,font=\footnotesize,minimum height=10mm] (loss) at (sm |- y) {損失 $-\log p_2$\\$=1.41$};
  \node[vec,minimum width=19mm] (g) at (z |- y) {勾配 $\vp-\vy$\\$+0.665$\\$-0.755$\\$+0.090$};
  \draw[arr] (x) -- (lin); \draw[arr] (lin) -- (z); \draw[arr] (z) -- (sm); \draw[arr] (sm) -- (p);
  \draw[arr,rounded corners=3pt] ($(p.south)+(-3mm,0)$) -- ++(0,-4.5mm) -| (loss.north);
  \draw[arr] (y) -- (loss); \draw[darr] (loss) -- (g);
  \draw[darr] (g.north) -- (z.south);
  \node[lbl,anchor=north] at ($(g.south)+(0,-0.8mm)$) {損失をロジットで微分した値};
  \node[lbl,anchor=north] at ($(y.south)+(0,-0.8mm)$) {データから};
\end{tikzpicture}
\caption{softmax 回帰の計算の順序と数値例。入力 $\vx$ に重み行列 $\mW$ を掛け、バイアス $\vb$ を足してロジット $\vz$ を作り、softmax で確率 $\vp$ に直す。
損失は、モデルが出した $\vp$ と、データから与えられる正解 $\vy$ の二つから計算する。
損失をロジットで微分すると、$\vp-\vy$ という簡単な形になる（破線）。
学習で動かすのはパラメータ $\mW$ と $\vb$ で、その勾配は $\vp-\vy$ から求まる（\eqref{eq:ml-softmax-param-grad}）。}
\label{fig:ml-softmax-flow}
\end{figure}

\subsection*{言語モデルの損失は、これと同じもの}

ここで、\chapref{ch:intro}の言語モデルを思い出してください。
「どうも仕方が」の次のトークンとして、モデルは「ない。」に確率 $0.559$ を割り当てました。
これは、$K=4096$ クラスの分類です。
実際の文章で次が「ない。」だったなら、損失は $-\log 0.559 = 0.58$ です。
次が「ないから」（確率 $0.121$）だったなら、損失は $-\log 0.121 = 2.11$ です。
何も学習していないモデルは、どのトークンにも、ほぼ同じ確率 $1/4096$ を割り当てます。
このときの損失は $-\log\frac{1}{4096} = \log 4096 = 8.32$ です。
実際、基盤モデルの学習は、ほぼこの値の $8.34$ から始まり、最後は $3.11$ まで下がりました（\chapref{ch:math}）。
言語モデルでは、文章の中の 1 か所（そこまでのトークンの並びと、次に来るトークンの組）が、データ 1 件にあたります。
言語モデルの学習とは、この交差エントロピー損失を、文章のあらゆる位置について平均し、小さくしていくことです。
詳しくは\chapref{ch:lm}で扱います。

\subsection*{勾配は「予測 $-$ 正解」になる}
\skimread{定理\ref{thm:ml-softmax-grad}の式 $\pd{\ell}{\vz}=\vp-\vy$ と、「数値例では」で始まる段落だけ押さえれば、先へ進めます。
導出は、\chapref{ch:nn}で誤差逆伝播法を学ぶときに戻れば足ります。}

勾配降下法を使うために、損失をロジットで微分します。

\begin{theorem}[softmax と交差エントロピーの勾配]\label{thm:ml-softmax-grad}
交差エントロピー損失 $\ell$ を、ロジット $z_k$ で偏微分すると
\begin{equation}
\pd{\ell}{z_k} = p_k - y_k
\qquad\text{（ベクトルで書けば }\pd{\ell}{\vz} = \vp - \vy\text{）}
\label{eq:ml-softmax-grad}
\end{equation}
となる。
\end{theorem}

導出します。
まず、損失をロジットで書き直します。
\eqref{eq:ml-softmax}の対数を取ると、商の対数は対数の差なので、$\log p_c = z_c - \log\sum_j e^{z_j}$ です。
損失は、次の二つの項に分かれます。
\begin{equation}
\ell = -\log p_c = -z_c + \log\sum_{j=1}^{K} e^{z_j}
\end{equation}
第 1 項 $-z_c$ を $z_k$ で微分します。
$k=c$ なら $-1$、そうでなければ $0$ です。これは $-y_k$ と書けます。
第 2 項は、「和 $u=\sum_j e^{z_j}$ を作る」と「$\log u$ を取る」の合成です。
$\log u$ の微分は $1/u$ です。
$u$ を $z_k$ で微分すると、和の中で $z_k$ を含む項だけが残り、$e^{z_k}$ です。
掛け合わせると $e^{z_k}/\sum_j e^{z_j}$ で、$p_k$ そのものです。
二つを合わせると、$\pd{\ell}{z_k} = p_k - y_k$ です。

数値例では $\vp-\vy = (0.665-0,\ 0.245-1,\ 0.090-0) = (0.665, -0.755, 0.090)$ です。
動かすのはパラメータ $\mW$ と $\vb$ ですが、この 1 件で更新すると、ロジットは勾配の符号と逆の向きに動きます。
正解のクラス 2 のロジットは上がり、ほかのクラスのロジットは下がります。
下げ幅は、間違って割り当てた確率に比例します。
予測 $\vp$ が正解 $\vy$ に近づくほど、勾配は 0 に近づきます。

パラメータについての勾配は、ここから 1 歩です。
スコアの式を成分で書くと、$z_k = \sum_j W_{kj}x_j + b_k$ です（この $j$ は、入力の成分の番号です）。
これを $W_{kj}$ で微分すると $x_j$、$b_k$ で微分すると 1 です。
$W_{kj}$ と $b_k$ が入っているロジットは $z_k$ だけなので、ほかのロジットは関係しません。
合成関数の微分により、次の式になります。
\begin{equation}
\pd{\ell}{W_{kj}} = (p_k - y_k)\,x_j,\qquad \pd{\ell}{b_k} = p_k - y_k
\label{eq:ml-softmax-param-grad}
\end{equation}
またしても「誤差 $\times$ 入力」です。
ロジスティック回帰の\eqref{eq:ml-logistic-grad}と、同じ形をしています。

\subsection*{実装}
\skimread{コードを 1 行ずつ追うのは、2 周目で構いません。押さえることは一つです。
PyTorch の \code{F.cross_entropy} は、ロジットと正解の番号を受け取り、softmax と交差エントロピー損失をまとめて計算します。}

平面上の点を 3 クラスに分けるモデルを、PyTorch で実装します。
データは 300 件で、入力 \code{X} の形は $(300, 2)$、正解 \code{y} はクラスの番号を並べた形 $(300,)$ のテンソルです。
数式ではクラスを 1 から、コードでは 0 から数えます。数式のクラス 2 は、コードの \code{1} です。

\snippet{softmax 回帰のモデルと損失}{code:ml-softmax-model}{ch03/softmax_regression.py}{model}

3 行目が\eqref{eq:ml-logits}です。
コードでは 1 データを 1 行に並べるので、\code{X @ W.T + b} と転置が付きます（\chapref{ch:pytorch}）。結果の形は $(N, K)$ です。
8 行目は、$\log p_k = z_k - \log\sum_j e^{z_j}$ を全データ・全クラスについて計算します。
\code{logsumexp} は、この対数と和と指数を、桁あふれを避けながら一度に計算する関数です（\chapref{ch:math}）。
9 行目の \code{log_p[torch.arange(len(y)), y]} は、各行から正解クラスの列だけを取り出します。
たとえば \code{y} が \code{[1, 0, 2]} なら、0 行目の 1 列目、1 行目の 0 列目、2 行目の 2 列目の 3 つを取り出します。
マイナスを付けて平均すれば、\eqref{eq:ml-cross-entropy}の平均です。

\snippet{手で導いた勾配 $\vp-\vy$}{code:ml-softmax-grad}{ch03/softmax_regression.py}{manual-grad}

3 行目で、データの件数 $N$ とクラスの数 $K$ を取り出します。
5 行目の \code{F.one_hot} が、クラスの番号を one-hot 表現に直します。結果の \code{Y} の形は $(N, K)$ です。
6 行目が\eqref{eq:ml-softmax-grad}で、\code{dz} の $i$ 行目が $i$ 件目の $\vp-\vy$ です。
損失は $N$ 件の平均なので、$N$ で割っています。
7 行目の \code{dz.T @ X} は、形 $(K, N)$ と $(N, 2)$ の行列積で、結果は $\mW$ と同じ形 $(K, 2)$ です。
その $(k, j)$ 成分は、\code{dz} の $k$ 列目と \code{X} の $j$ 列目の内積です。
つまり、\eqref{eq:ml-softmax-param-grad}の $(p_k-y_k)\,x_j$ を、$N$ 件について足したものです。
\code{dz.sum(dim=0)} も同じように、$p_k-y_k$ を $N$ 件について足したもので、$\vb$ の勾配になります。

% [figpass] 案 2
\figref{fig:ml-dz-shape}は、7 行目の行列積の形を描いたものです。
\code{dz.T} の行（クラス）と \code{X} の列（入力の成分）を一つずつ選び、データの番号の向きに $N$ 個の積を足すので、$N$ の軸は消え、$(K,2)$ の形が残ります。

\begin{figure}[tbp]
\centering
\begin{tikzpicture}[x=1mm,y=1mm,
  hd/.style={font=\footnotesize,align=center},
  sm/.style={font=\scriptsize,text=bookgray,align=center,inner sep=1pt}]
  % dz.T: (K, N) = (3, 300)。横の長さは縮めて描く
  \fill[figorange] (0,-3) rectangle (40,-6);
  \draw[black!60] (0,0) rectangle (40,-9);
  \draw[black!40] (0,-3) -- (40,-3); \draw[black!40] (0,-6) -- (40,-6);
  \draw[bookorange,thick] (0,-3) rectangle (40,-6);
  \foreach \r in {1,2,3} \node[sm,anchor=east] at (-0.5,-3*\r+1.5) {$k=\r$};
  \draw[decorate,decoration={brace,amplitude=3pt}] (0,1) -- node[above=2pt,sm]{$N=300$（データの番号）} (40,1);
  \node[hd,anchor=north] at (20,-10.5) {\code{dz.T}\\形 $(K,N)=(3,300)$\\$i$ 列目は $i$ 件目の $(\vp-\vy)/N$};
  \node[font=\small] at (45,-4.5) {$@$};
  % X: (N, 2)。縦の長さは縮めて描く
  \fill[figblue] (50,10.5) rectangle (58,-19.5);
  \draw[black!60] (50,10.5) rectangle (66,-19.5);
  \draw[black!40] (58,10.5) -- (58,-19.5);
  \draw[bookblue,thick] (50,10.5) rectangle (58,-19.5);
  \node[sm,anchor=south] at (54,10.7) {$j=1$}; \node[sm,anchor=south] at (62,10.7) {$j=2$};
  \draw[decorate,decoration={brace,amplitude=3pt}] (67,10.5) -- node[right=5pt,sm]{$N=300$} (67,-19.5);
  \node[hd,anchor=north] at (58,-20.5) {\code{X}\\形 $(N,2)=(300,2)$};
  \node[font=\small] at (88,-4.5) {$=$};
  % 結果: (K, 2)
  \fill[figpurple] (93,-3) rectangle (99,-6);
  \draw[black!60] (93,0) rectangle (105,-9);
  \draw[black!40] (93,-3) -- (105,-3); \draw[black!40] (93,-6) -- (105,-6); \draw[black!40] (99,0) -- (99,-9);
  \draw[bookpurple,thick] (93,-3) rectangle (99,-6);
  \node[hd,anchor=north] at (99,-10.5) {\code{dz.T @ X}\\形 $(K,2)=(3,2)$\\$\mW$ の勾配};
  % 説明
  \node[sm,anchor=north] at (48,-30) {橙の行と青の列の内積（$N$ 個の積の和）が、紫のマス 1 つになる。$N$ の軸は消え、$K$ と 2 の軸が残る。};
\end{tikzpicture}
\caption{\coderef{code:ml-softmax-grad}の 7 行目 \code{dz.T @ X} の形（模式図。$N$ の向きの長さは縮めて描いた）。
紫のマスは、$\pd{\calL}{W_{21}}=\frac{1}{N}\sum_{i}(p_2^{(i)}-y_2^{(i)})\,x_1^{(i)}$ で、\eqref{eq:ml-softmax-param-grad}を $N$ 件について平均したもの。
ほかのマスも同じで、行 $k$ と列 $j$ の組ごとに 1 つずつ計算される。}
\label{fig:ml-dz-shape}
\end{figure}
% [/figpass]

手で導いたこの勾配を、同じ損失から自動微分で求めた勾配と比べました。
乱数で決めたパラメータで比べると、差は最大でも $1.2\times10^{-7}$ で、計算の丸め誤差の範囲でした。導出は正しいことが確かめられました。
なお、PyTorch の \code{F.cross_entropy} は、ロジットと正解の番号を受け取り、softmax と交差エントロピーをまとめて計算します。
同じ入力で、自作の関数と同じ値（$7.27554$）を返しました。今後は、この関数を使います。

\subsection*{結果：境界はまっすぐになる}

3 つの塊のデータ（300 件）で、softmax 回帰を学習させます。
学習のループは、\coderef{code:ml-torch-loop}と同じ 4 手順です（\file{code/ch03/softmax_regression.py} の \code{train} 関数）。
パラメータをすべて 0 から始め、学習率 $0.5$ で 300 ステップ学習しました。
最初の損失は $1.0986$ で、これは $\log 3$ です。
パラメータが 0 なら、3 クラスに同じ確率 $1/3$ を割り当てるからです。
\figref{fig:ml-blobs}の左が、学習中の損失の動きです。
損失は最初の 10 ステップで $0.23$ まで急に下がり、そのあとはゆっくり下がって、300 ステップ後に $0.170$ になりました。
確率が最大のクラスを予測とし、それが正解と一致した割合を\textbf{正解率}と呼びます。
正解率は 93.7\% で、同じ作り方で新しく作った、学習に使っていない 300 点でも 94.3\% です。

\begin{figure}[tbp]
\centering
\pgfplotsset{mlthree/.style={width=0.385\linewidth,height=0.31\linewidth,
  title style={font=\footnotesize,yshift=-1.5mm},tick label style={font=\scriptsize},
  label style={font=\scriptsize}}}
\begin{tikzpicture}
\begin{axis}[mlthree,title={学習中の損失と正解率},xlabel={ステップ},
  xmin=0,xmax=100,ymin=0,ymax=1.2]
\addplot[bookaccent,very thick] table[col sep=comma,x=step,y=loss]{code/ch03/out/softmax_curve.csv};
\addplot[bookblue,thick,densely dashed] table[col sep=comma,x=step,y=acc]{code/ch03/out/softmax_curve.csv};
\node[font=\scriptsize,text=bookblue,anchor=north east] at (axis cs:97,0.91) {正解率};
\node[font=\scriptsize,text=bookaccent,anchor=south east] at (axis cs:97,0.2) {損失};
\end{axis}
\end{tikzpicture}\hfill
\begin{tikzpicture}
\begin{axis}[mlthree,xlabel={$x_1$},ylabel={$x_2$},ylabel style={yshift=-2mm},
  xmin=-5,xmax=5,ymin=-4,ymax=5,grid=none,title={塊：正解率 93.7\%},
  colormap={cls}{color(0)=(figblue) color(1)=(figorange) color(2)=(figgreen)},point meta min=0,point meta max=2]
\addplot[matrix plot*,mesh/cols=61,point meta=explicit,shader=flat corner]
  table[col sep=comma,x=x1,y=x2,meta=cls]{code/ch03/out/blobs_grid.csv};
\addplot[only marks,mark=*,mark size=0.9pt,bookblue,restrict expr to domain={\thisrow{cls}}{0:0}]
  table[col sep=comma,x=x1,y=x2]{code/ch03/out/blobs_points.csv};
\addplot[only marks,mark=square*,mark size=0.9pt,bookorange,restrict expr to domain={\thisrow{cls}}{1:1}]
  table[col sep=comma,x=x1,y=x2]{code/ch03/out/blobs_points.csv};
\addplot[only marks,mark=triangle*,mark size=1.2pt,bookgreen,restrict expr to domain={\thisrow{cls}}{2:2}]
  table[col sep=comma,x=x1,y=x2]{code/ch03/out/blobs_points.csv};
\end{axis}
\end{tikzpicture}\hfill
\begin{tikzpicture}
\begin{axis}[mlthree,xlabel={$x_1$},
  xmin=-1.2,xmax=1.2,ymin=-1.2,ymax=1.2,grid=none,title={渦巻き：正解率 52.7\%},
  colormap={cls}{color(0)=(figblue) color(1)=(figorange) color(2)=(figgreen)},point meta min=0,point meta max=2]
\addplot[matrix plot*,mesh/cols=61,point meta=explicit,shader=flat corner]
  table[col sep=comma,x=x1,y=x2,meta=cls]{code/ch03/out/spiral_grid.csv};
\addplot[only marks,mark=*,mark size=0.9pt,bookblue,restrict expr to domain={\thisrow{cls}}{0:0}]
  table[col sep=comma,x=x1,y=x2]{code/ch03/out/spiral_points.csv};
\addplot[only marks,mark=square*,mark size=0.9pt,bookorange,restrict expr to domain={\thisrow{cls}}{1:1}]
  table[col sep=comma,x=x1,y=x2]{code/ch03/out/spiral_points.csv};
\addplot[only marks,mark=triangle*,mark size=1.2pt,bookgreen,restrict expr to domain={\thisrow{cls}}{2:2}]
  table[col sep=comma,x=x1,y=x2]{code/ch03/out/spiral_points.csv};
\end{axis}
\end{tikzpicture}
\caption{左：塊のデータでの、最初の 100 ステップの損失（赤）と正解率（青の破線。0〜1 の割合）。縦軸は二つに共通。
中央と右：softmax 回帰が学習した領域分け（横軸 $x_1$、縦軸 $x_2$）。点がデータで、クラス 1 が青の丸、クラス 2 が橙の四角、クラス 3 が緑の三角。
背景の色（点と同じ系統の薄い色）は、モデルが予測するクラス（細かいマス目ごとに塗ったので境界が階段状に見えるが、実際の境界は直線）。
渦を巻くデータ（右）は、直線の境界では分けられない。}
\label{fig:ml-blobs}
\end{figure}

\figref{fig:ml-blobs}の中央が、学習したモデルの領域分けです。
背景の色は、その場所をモデルがどのクラスと予測するかを表します。
3 つの塊が、おおよそ正しく塗り分けられています。
塊どうしが重なっている場所では、違うクラスの点が入り混じっています。
直線の境界をどう引いても、一部は間違えます。正解率が 100\% にならないのは、そのためです。

境界が\textbf{まっすぐ}であることに注目してください。
クラス $k$ とクラス $l$ の境界は、二つのスコアが等しくなる場所 $z_k = z_l$ です。
$z_k$ も $z_l$ も $x_1, x_2$ の 1 次式なので、この条件は直線の式になります。
入力 $x_1, x_2$ をそのまま使うかぎり、softmax 回帰は、どれだけ学習しても、直線の境界しか作れません。
\figref{fig:ml-blobs}の右は、別のデータに、同じモデルを学習させた結果です。
このデータでは、3 つのクラスの点が、それぞれ中心から渦巻きの形に伸びる帯の上に並んでいます。
正解率は 52.7\% にとどまります。
3 クラスなので、でたらめに答えても約 33\% は当たります。
そこから 20 ポイントほどしか上がらず、半分近くを間違えています。
この限界を越えることが、次章の出発点になります。

最後に、この章の三つのモデルを\tabref{tab:ml-reg-vs-cls}に並べます。
どれも、正解についての確率分布を出力するようにモデルを作り、負の対数尤度を損失にしたものです。
損失を、線形回帰では予測 $\hat{y}$ で、分類ではロジットで微分すると、どれも「予測 $-$ 正解」の形になります。

\begin{table}[tbp]
\centering
\caption{この章の三つのモデル。$\hat{y}$ は線形回帰の予測。
線形回帰の損失は、負の対数尤度（\eqref{eq:ml-gauss-nll}）から、パラメータに関係しない項と定数倍を除いたもので、微分には 2 乗から来る 2 倍が付く。}
\label{tab:ml-reg-vs-cls}
\small
\begin{tabularx}{\linewidth}{l>{\hsize=.9\hsize}X>{\hsize=.95\hsize}X>{\hsize=1.15\hsize}X}
\toprule
 & 線形回帰 & ロジスティック回帰 & softmax 回帰 \\
\midrule
正解 $y$ & 数 & 0 か 1 & $K$ 種類のどれか \\
モデルの出力 & $\hat{y}=\vw^\T\vx+b$ & $p=\sigmoid(\vw^\T\vx+b)$ & \mbox{$\vp=\softmax(\mW\vx+\vb)$} \\
人が選ぶ $y$ の分布 & $\hat{y}$ を中心とする正規分布 & 確率 $p$ で 1、$1-p$ で 0 & クラス $k$ が確率 $p_k$ \\
損失 $\ell$ & $(\hat{y}-y)^2$ & \eqref{eq:ml-bce}（2 クラスの交差エントロピー損失） & $-\log p_c$（交差エントロピー損失） \\
$\hat{y}$ かロジットでの微分 & $\pd{\ell}{\hat{y}}=2(\hat{y}-y)$ & $\pd{\ell}{z}=p-y$ & $\pd{\ell}{\vz}=\vp-\vy$ \\
\bottomrule
\end{tabularx}
\end{table}

この節では、種類を答える問題を扱いました。
正解のクラスの確率の対数にマイナスを付けたものが交差エントロピー損失で、言語モデルの損失と同じものです。

\section{確率的勾配降下法とミニバッチ}\label{sec:ml-sgd}

ここまでは、1 ステップごとに全データで勾配を計算してきました。
データが 100 件なら問題ありません。
しかし、基盤モデルの学習データは約 250 万トークンあり、実際の LLM では兆の単位になります。
この節の問いは、「データが膨大なとき、1 ステップをどう軽くするか」です。
4 要素のうち、こんどは最適化に手を入れます。

\textbf{本来計算したいもの}は、\textbf{全データの平均の勾配}です。
損失はデータ 1 件ごとの損失の平均なので（定義\ref{def:ml-erm}）、勾配も 1 件ごとの勾配の平均です（和の微分は、項ごとの微分の和だからです）。
\textbf{それが難しい理由}は、手間です。
これまでの勾配降下法は、1 回の更新のたびに、$N$ 件すべての勾配を計算して平均していました。
データが 1 万件なら、1 万件分の計算をして、やっと 1 回の更新です。

\textbf{代わりにすること}は、一部のデータでの代用です。
平均は、無作為に選んだ一部からでも見積もれます。
世論調査が、全国民には尋ねず、無作為に選んだ千人ほどの答えから全体の傾向を見積もるのと同じ考え方です（全国民が全データ、千人がその一部にあたります）。
そこで、全データの平均の勾配の代わりに、無作為に選んだ一部のデータの平均の勾配を使います。
選んだ一部のデータを\term{ミニバッチ}{minibatch}、その件数 $B$ を\term{バッチサイズ}{batch size}と呼びます。

\textbf{それでよい理由}は、一部の平均の勾配が、全データの平均の勾配の、偏りのない見積もりだからです。
これを実験で確かめます。
文字数とトークン数のデータを 1 万件に増やします。
$w=b=0$ の地点で、1 万件の中から $B$ 件を無作為に選び、その $B$ 件だけで勾配 $\pd{\calL}{w}$ を計算することを、2000 回繰り返しました。
結果が次です。「平均」と「標準偏差」は、2000 回分の勾配について求めたものです。

\outputfile{ch03/out/sgd_noise.txt}

全データで計算した勾配は $-1.777$ です。
$B=1$ のとき、2000 回の勾配の平均は $-1.776$、標準偏差は $1.446$ でした。
$B=100$ では、平均が $-1.770$、標準偏差が $0.139$ です。
平均は、どちらも全データの勾配とほぼ同じです。
ミニバッチの勾配は、1 回ごとには全データの勾配から外れますが、何回分も平均すれば、全データの勾配に近づいていきます。
これが、偏りのない見積もり（統計の言葉では「不偏」）という意味です。
一方、標準偏差は見積もりのばらつきを表します。
$B$ を 1、10、100、1000 と増やすと、標準偏差は $1.446$、$0.459$、$0.139$、$0.045$ と減りました。
出力の値で、隣どうしの比を取ります。
$1.4460\div0.4586=3.15$、$0.4586\div0.1393=3.29$、$0.1393\div0.0445=3.13$ です。
$B$ を 10 倍にするたびに約 $1/3.2$ 倍、100 倍にすると約 $1/10$ 倍です。
$\sqrt{10}\approx 3.2$ なので、ばらつきは $1/\sqrt{B}$ に比例して減っています。
ここに、ミニバッチを使う理由があります。
計算量を 100 倍かけても、勾配の精度は 10 倍にしかなりません。
それなら、粗い見積もりで何度も進むほうが得です。
偏りがないので、更新を重ねれば、外れは打ち消し合います。

% [figpass] 案 6
\figref{fig:ml-grad-noise}は、この実験を、\file{code/ch03/sweep_grad_noise.py} で $B$ を 10 通りに増やして走らせ直したものです（地点 $w=b=0$ と、2000 回選び直すことは同じ）。
上の 3 枚は、2000 個の勾配のヒストグラムです。
どの $B$ でも、2000 個の勾配の平均は全データの勾配（赤い縦線）の位置にあり、$B$ を大きくするほど山が細くなって、縦線の近くに集まります。
下の図では、標準偏差の点が、$1/\sqrt{B}$ に比例する直線にほぼ乗っています。

\begin{figure}[tbp]
\centering
\pgfplotsset{mlhist/.style={width=0.37\linewidth,height=0.25\linewidth,
  xmin=-6,xmax=0.5,ymin=0,xtick={-6,-4,-2,0},
  tick label style={font=\scriptsize},label style={font=\scriptsize},
  title style={font=\footnotesize,yshift=-1.5mm},xlabel={勾配 $\pd{\calL}{w}$}}}
\begin{tikzpicture}
\begin{axis}[mlhist,title={$B=1$},ylabel={確率密度},ymax=0.8]
\addplot[ybar interval,fill=figblue,draw=bookblue,thin] table[col sep=comma,x expr=\thisrow{center}-0.05,y=b1]{code/ch03/out/sweep_grad_hist.csv};
\draw[bookaccent,thick] (axis cs:-1.7766,0) -- (axis cs:-1.7766,0.8);
\end{axis}
\end{tikzpicture}\hfill
\begin{tikzpicture}
\begin{axis}[mlhist,title={$B=10$},ymax=0.9]
\addplot[ybar interval,fill=figblue,draw=bookblue,thin] table[col sep=comma,x expr=\thisrow{center}-0.05,y=b10]{code/ch03/out/sweep_grad_hist.csv};
\draw[bookaccent,thick] (axis cs:-1.7766,0) -- (axis cs:-1.7766,0.9);
\end{axis}
\end{tikzpicture}\hfill
\begin{tikzpicture}
\begin{axis}[mlhist,title={$B=100$},ymax=3.2]
\addplot[ybar interval,fill=figblue,draw=bookblue,thin] table[col sep=comma,x expr=\thisrow{center}-0.05,y=b100]{code/ch03/out/sweep_grad_hist.csv};
\draw[bookaccent,thick] (axis cs:-1.7766,0) -- (axis cs:-1.7766,3.2);
\end{axis}
\end{tikzpicture}

\vspace{1mm}
\begin{tikzpicture}
\begin{loglogaxis}[width=0.6\linewidth,height=0.3\linewidth,xlabel={バッチサイズ $B$},ylabel={勾配の標準偏差},
  xmin=0.7,xmax=1500,ymin=0.03,ymax=2.5,
  tick label style={font=\scriptsize},label style={font=\footnotesize},
  legend style={font=\scriptsize,at={(0.98,0.97)},anchor=north east,draw=black!30},legend cell align=left]
\addplot[black!50,thick,densely dashed] table[col sep=comma,x=B,y=ref]{code/ch03/out/sweep_grad_noise.csv};
\addplot[only marks,mark=*,mark size=1.6pt,bookblue] table[col sep=comma,x=B,y=std]{code/ch03/out/sweep_grad_noise.csv};
\legend{$1/\sqrt{B}$ に比例する線,実測（2000 回）}
\end{loglogaxis}
\end{tikzpicture}
\caption{ミニバッチの勾配のばらつき（$w=b=0$ の地点、1 万件から $B$ 件を選ぶことを 2000 回）。
上：2000 個の勾配のヒストグラム。赤い縦線は全データの勾配 $-1.777$。縦軸の範囲は 3 枚で違う（$B$ が大きいほど山が細く高い）。
$B=1$ の山が左右対称でないのは、1 件の勾配 $-2x^{(i)}y^{(i)}$ が、長い断片ほど大きく負になるため。
下：標準偏差（両対数目盛り）。破線は、$B=1$ の値を $\sqrt{B}$ で割った線。
乱数の選び方が\file{code/ch03/sgd_minibatch.py}と違うので、$B=10$ などの値は、本文の出力と小数第 2 位で違う。}
\label{fig:ml-grad-noise}
\end{figure}
% [/figpass]

では、これは、データが大量にあるからこそ使える方法なのでしょうか。
そうではありません。
見積もりに偏りがないことは、データの件数 $N$ に関係なく成り立つので、データが 100 件でも同じように使えます。
ただし、データが少なければ、全データで毎回計算しても手間は小さく、利点も小さくなります。
データが大量にあるときは、この方法を使わなければ、学習を進められません。

以上をまとめて、1 回の更新（1 ステップ）ごとに $B$ 件のミニバッチ $\mathcal{B}$ を取り出し、その平均の勾配でパラメータを更新します。
\begin{equation}
\vtheta \leftarrow \vtheta - \eta\,\nabla_{\vtheta}\Bigl(\frac{1}{B}\sum_{i\in\mathcal{B}} \ell\bigl(f_\theta(\vx^{(i)}),\, y^{(i)}\bigr)\Bigr)
\label{eq:ml-sgd}
\end{equation}
$\mathcal{B}$ は、そのステップで使うミニバッチ（取り出した $B$ 件のデータの番号の集まり）で、件数の $B$ とは別の記号です。
入力を太字の $\vx^{(i)}$ と書いたのは、入力が複数の数を並べたベクトルの場合も含めるためです。
定義\ref{def:ml-erm}の全データの平均が、ミニバッチの中だけの平均 $\frac{1}{B}\sum_{i\in\mathcal{B}}$ に置き換わっています。
この方法を\term{確率的勾配降下法}{stochastic gradient descent, SGD}と呼びます。

「確率的」とは、データを無作為に選ぶので、進む向きが毎回偶然に左右される、という意味です。
この考え方は 1951 年の論文\cite{robbins1951stochastic}にさかのぼります。

大事なのは、\textbf{ミニバッチをステップごとに取り替える}ことです。
同じ $B$ 件だけを使い続けると、その $B$ 件だけに合ったパラメータになってしまいます。

実際の学習では、ミニバッチを次の手順で作るのが普通です（\figref{fig:ml-minibatch-epoch}）。
(1) 全データの順番を無作為に混ぜ、(2) 先頭から $B$ 件ずつに切り分け、(3) 切り分けた組を順に 1 つずつ使って、1 回ずつパラメータを更新します。
すべての組を使い終えると、全データをちょうど 1 回ずつ使ったことになります。
この一巡を 1 \term{エポック}{epoch}と呼び、次のエポックでは、順番を混ぜ直して同じことを繰り返します。
つまり、データ全体を $B$ 件ずつの組に分け、1 回の更新では、全データの平均の代わりに、そのうち 1 組の平均を使います。
組の一つ一つがミニバッチで、1 組の件数がバッチサイズです。

\begin{figure}[tbp]
\centering
\begin{tikzpicture}[x=6.2mm,y=5.6mm,
  c/.style={cell,minimum width=6.2mm,minimum height=5.6mm,inner sep=0pt},
  note/.style={font=\footnotesize,align=center,inner sep=1pt}]
  % 1 段目: 全データ 12 件（番号順）
  \foreach \i in {1,...,12} \node[c] at (\i,0) {\i};
  \node[lbl,anchor=east] at (0.3,0) {全データ（$N=12$ 件）};
  \draw[arr] (6.5,-0.55) -- node[right,lbl]{(1) 順番を無作為に混ぜる} (6.5,-1.4);
  % 2 段目: 混ぜたあと。4 件ずつ色分け
  \foreach \v [count=\i from 1] in {7,2,11,5} \node[c,fill=figblue] at (\i,-2) {\v};
  \foreach \v [count=\i from 5] in {9,1,12,4} \node[c,fill=figorange] at (\i,-2) {\v};
  \foreach \v [count=\i from 9] in {3,10,6,8} \node[c,fill=figgreen] at (\i,-2) {\v};
  \node[lbl,anchor=east,align=right] at (0.3,-2) {(2) $B=4$ 件ずつに\\切り分ける};
  % 切れ目
  \draw[very thick,black!70] (4.5,-1.4) -- (4.5,-2.6);
  \draw[very thick,black!70] (8.5,-1.4) -- (8.5,-2.6);
  % 3 段目: 各ミニバッチで 1 回ずつ更新
  \draw[arr] (2.5,-2.6) -- (2.5,-3.3);
  \draw[arr] (6.5,-2.6) -- (6.5,-3.3);
  \draw[arr] (10.5,-2.6) -- (10.5,-3.3);
  \node[note,anchor=north] at (2.5,-3.35) {1 回目の更新};
  \node[note,anchor=north] at (6.5,-3.35) {2 回目の更新};
  \node[note,anchor=north] at (10.5,-3.35) {3 回目の更新};
  \node[lbl,anchor=east,align=right] at (0.3,-3.55) {(3) 1 組ずつ、その平均の\\勾配で更新する};
\end{tikzpicture}
\caption{ミニバッチの作り方と 1 エポック（模式図。データ 12 件、バッチサイズ 4 の例で、マスの中の数はデータの番号）。
順番を混ぜて 4 件ずつに切り分けると、3 つのミニバッチができ（青、橙、緑）、1 回の更新にはその 1 つだけを使う。
3 回更新すると、12 件すべてをちょうど 1 回ずつ使ったことになる（1 エポック）。次のエポックでは、(1) の混ぜ直しからやり直す。}
\label{fig:ml-minibatch-epoch}
\end{figure}

データが $N$ 件なら、1 エポックは $N/B$ ステップです（割り切れなければ切り上げ）。
1 万件を $B=10$ で一巡するなら、$10000\div10=1000$ ステップです。
先ほどの実験では、毎回 $B$ 件を無作為に選び直しましたが、次の学習のコードは、「混ぜて、切り分けて、順に使う」手順です。

\snippet{ミニバッチ SGD のループ}{code:ml-sgd-loop}{ch03/sgd_minibatch.py}{sgd-loop}

2 行目で、0 から $N-1$ までのデータの番号を、無作為に並べ替えます。
3〜4 行目で、その先頭から $B$ 個ずつ番号を取り出します。これがミニバッチです。
5 行目は、\coderef{code:ml-linreg-model}の \code{grad} に、選んだ $B$ 件だけを渡します。
全データを渡していた\coderef{code:ml-gd}との違いは、ここだけです。6〜7 行目の更新は変わりません。
3 行目からの内側のループを終えると 1 エポックで、1 行目の外側のループがそれを繰り返します。
\code{rng} は乱数を作る NumPy の部品、\code{n_epochs} はエポックの数、\code{eta} は学習率 $\eta$ です。
2 行目が外側のループの中にあるので、順番はエポックごとに混ぜ直されます。

では、同じ計算量でどこまで進めるかを比べます。
学習率を $0.1$ にそろえ、バッチサイズを変えて学習しました。
ミニバッチを使う学習は、どのバッチサイズでも 3 エポック（3 万件の処理）です。$B=10$ なら 3,000 ステップにあたります。
公平に比べるため、横軸をステップ数ではなく、\textbf{処理したデータの件数}にします（\figref{fig:ml-sgd}の左）。
全データを使う勾配降下法は、1 万件を処理して、やっと 1 ステップです。
1 万件を処理した時点の損失は $0.167$ でした。
$B=10$ の SGD は、同じ 1 万件で 1000 ステップ進み、損失は $0.0093$ です。
これは、到達できる最小値 $0.0091$ のすぐ近くです。
勾配降下法がこの水準に達するには、100 万件（100 エポック）の処理が必要でした。

\begin{figure}[tbp]
\centering
\begin{tikzpicture}
\begin{loglogaxis}[width=0.52\linewidth,height=0.315\linewidth,xlabel={処理したデータの件数},ylabel={損失（1 万件で評価）},
  xmin=1,xmax=3e6,ymin=0.007,ymax=9,
  tick label style={font=\scriptsize},label style={font=\footnotesize},
  legend columns=2,
  legend style={font=\scriptsize,at={(0.98,0.97)},anchor=north east,draw=black!30,/tikz/every even column/.append style={column sep=1.5mm}},legend cell align=left]
\addplot[black,very thick] table[col sep=comma,x=seen,y=loss]{code/ch03/out/sgd_curve_full.csv};
\addplot[bookgreen,thick] table[col sep=comma,x=seen,y=loss]{code/ch03/out/sgd_curve_b100.csv};
\addplot[bookaccent,thick] table[col sep=comma,x=seen,y=loss]{code/ch03/out/sgd_curve_b10.csv};
\addplot[bookblue,thin] table[col sep=comma,x=seen,y=loss]{code/ch03/out/sgd_curve_b1.csv};
\legend{全データ,$B=100$,$B=10$,$B=1$}
\addplot[only marks,mark=*,mark size=1.5pt,black] coordinates {(10000,0.16717)};
\node[font=\scriptsize,anchor=west,align=left,inner sep=1.5pt] at (axis cs:15000,0.165) {全データの\\1 ステップ目};
\end{loglogaxis}
\end{tikzpicture}\hfill
\begin{tikzpicture}
\begin{axis}[width=0.48\linewidth,height=0.315\linewidth,xlabel={重み $w$},ylabel={バイアス $b$},
  xmin=-0.05,xmax=0.8,ymin=-0.12,ymax=0.3,
  tick label style={font=\scriptsize},label style={font=\footnotesize}]
\addplot[black!22,thin,empty line=jump] table[col sep=comma,x=w,y=b]{code/ch03/out/loss_contours_big.csv};
\addplot[bookblue,thin,restrict expr to domain={\thisrow{step}}{0:120}] table[col sep=comma,x=w,y=b]{code/ch03/out/sgd_path_b1.csv};
\addplot[bookaccent,thick] table[col sep=comma,x=w,y=b]{code/ch03/out/sgd_path_b10.csv};
\addplot[black,very thick] table[col sep=comma,x=w,y=b]{code/ch03/out/sgd_path_full.csv};
\addplot[only marks,mark=*,mark size=2.6pt,white] coordinates {(0.6033,-0.0046)};
\addplot[only marks,mark=x,mark size=3.5pt,very thick,black] coordinates {(0.6033,-0.0046)};
\addplot[only marks,mark=*,mark size=2pt,black,mark options={fill=white}] coordinates {(0.17766,0.13192)};
\node[font=\scriptsize,anchor=south west,align=left,inner sep=1.5pt] (ml-sgd-first) at (axis cs:-0.045,0.2) {黒の 1 ステップ目\\（1 万件を処理）};
\draw[black!60,thin] (axis cs:0.115,0.198) -- (axis cs:0.168,0.142);
\end{axis}
\end{tikzpicture}
\caption{左：処理したデータの件数と損失（両軸とも対数目盛り）。線の色は左右で共通で、黒が全データ、緑が $B=100$、赤が $B=10$、青が $B=1$。
右：パラメータ $(w,b)$ が動いた跡。$(0,0)$ から始めた最初の 300 ステップを描いた（$B=1$ は 120 ステップまで。$B=100$ は描いていない）。$\times$ 印が谷底。
右の図でそろえたのはステップ数で、処理したデータは、黒が 300 万件、赤が 3,000 件、青が 120 件である。
灰色の細い線は等高線で、\figref{fig:ml-landscape}と同じ細長い谷の一部を拡大している。}
\label{fig:ml-sgd}
\end{figure}

\figref{fig:ml-sgd}の右は、パラメータ $(w,b)$ が動いた跡（ステップごとの $(w,b)$ を順につないだ線）を、ステップ数をそろえて描いたものです。
全データの跡（黒）はなめらかです。
ミニバッチの跡（赤と青）は、勾配の見積もりがぶれるので、ふらつきます。
それでも、平均すれば正しい向きなので、同じ谷底へ向かいます。
赤は、黒の 1000 分の 1 のデータしか処理していないのに、谷底の近くまで来ています。
ただし、谷底に着いても止まらず、見積もりのぶれのせいで、まわりを揺れ続けます。
左の図の青い線（$B=1$）が、おおむね $0.01$ から $0.02$ の間で上下し、ときには $0.04$ 近くまで上がるのは、このためです。ただし、この線は記録を間引いて描いたもので、毎ステップの損失はもっと高く上がることがあります（この節の少し先で見ます）。
揺れは、バッチサイズを大きくするか、学習率を小さくすると、小さくなります。
% [figpass] 案 7
\figref{fig:ml-sgd-eta}は、これを確かめたものです。
\file{code/ch03/sweep_sgd_eta.py} で、$B$ を 3 通り、学習率を 0.003〜0.3 の 5 通りに変え、乱数の種を 3 つずつ走らせました。
右の図は、谷底から出発して 1 エポック走らせたときに、損失が最小値をどれだけ上回ったかで、谷底のまわりの揺れだけを表します。
どの $B$ でも、学習率を下げるほど揺れは小さく、同じ学習率なら $B$ が大きいほど小さくなっています。
左の図は、$(0,0)$ から 3 エポック（3 万件）走らせたあとの同じ量で、谷底まで進む速さと揺れの両方を含みます。
$B=100$ で学習率を下げすぎると、3 エポックでは谷底に着かず、かえって損失が大きく残ります。
揺れを抑えることと速く進むことは両立しないので、学習率と $B$ は組にして選びます。
\figref{fig:ml-sgd}の左の青い線は、記録を間引いて描いたものでした。
毎ステップで見ると、$B=1$・学習率 $0.1$ の損失は、3 エポック目の 1.4\% ほどのステップで 0.04 を超え、最大で 0.24 まで上がりました（\file{code/ch03/out/sweep_sgd_eta.txt}）。

\begin{figure}[tbp]
\centering
\pgfplotsset{mlsgdeta/.style={height=0.33\linewidth,
  xmin=0.002,xmax=0.45,xlabel={学習率 $\eta$},
  tick label style={font=\scriptsize},label style={font=\footnotesize},
  title style={font=\footnotesize,yshift=-1.5mm}}}
\begin{tikzpicture}
\begin{loglogaxis}[mlsgdeta,width=0.53\linewidth,title={$(0,0)$ から 3 エポック},
  ylabel={損失 $-$ 最小値},ymin=1e-5,ymax=1,
  legend to name=mlsgdetaleg,legend columns=3,legend style={font=\scriptsize,draw=black!30,/tikz/every even column/.append style={column sep=3mm}}]
\addplot[bookblue,thick,mark=*,mark size=1.2pt,error bars/.cd,y dir=both,y explicit] table[col sep=comma,x=eta,y=excess,y error plus expr=\thisrow{hi}-\thisrow{excess},y error minus expr=\thisrow{excess}-\thisrow{lo}]{code/ch03/out/sweep_sgd_from0_b1.csv};
\addplot[bookaccent,thick,mark=square*,mark size=1.2pt,error bars/.cd,y dir=both,y explicit] table[col sep=comma,x=eta,y=excess,y error plus expr=\thisrow{hi}-\thisrow{excess},y error minus expr=\thisrow{excess}-\thisrow{lo}]{code/ch03/out/sweep_sgd_from0_b10.csv};
\addplot[bookgreen,thick,mark=triangle*,mark size=1.6pt,error bars/.cd,y dir=both,y explicit] table[col sep=comma,x=eta,y=excess,y error plus expr=\thisrow{hi}-\thisrow{excess},y error minus expr=\thisrow{excess}-\thisrow{lo}]{code/ch03/out/sweep_sgd_from0_b100.csv};
\legend{$B=1$,$B=10$,$B=100$}
\end{loglogaxis}
\end{tikzpicture}\hfill
\begin{tikzpicture}
\begin{loglogaxis}[mlsgdeta,width=0.49\linewidth,title={谷底から 1 エポック（揺れ）},
  ymin=3e-8,ymax=1]
\addplot[bookblue,thick,mark=*,mark size=1.2pt,error bars/.cd,y dir=both,y explicit] table[col sep=comma,x=eta,y=excess,y error plus expr=\thisrow{hi}-\thisrow{excess},y error minus expr=\thisrow{excess}-\thisrow{lo}]{code/ch03/out/sweep_sgd_frombest_b1.csv};
\addplot[bookaccent,thick,mark=square*,mark size=1.2pt,error bars/.cd,y dir=both,y explicit] table[col sep=comma,x=eta,y=excess,y error plus expr=\thisrow{hi}-\thisrow{excess},y error minus expr=\thisrow{excess}-\thisrow{lo}]{code/ch03/out/sweep_sgd_frombest_b10.csv};
\addplot[bookgreen,thick,mark=triangle*,mark size=1.6pt,error bars/.cd,y dir=both,y explicit] table[col sep=comma,x=eta,y=excess,y error plus expr=\thisrow{hi}-\thisrow{excess},y error minus expr=\thisrow{excess}-\thisrow{lo}]{code/ch03/out/sweep_sgd_frombest_b100.csv};
\end{loglogaxis}
\end{tikzpicture}

\vspace{1mm}
\ref{mlsgdetaleg}
\caption{ミニバッチ SGD で、バッチサイズ $B$ と学習率 $\eta$ を変えたときの、損失が最小値 $0.00909$ を上回る量（最後のエポックの毎ステップの損失の平均から最小値を引いた値。両軸とも対数目盛り）。
データ（1 万件）、ミニバッチの作り方、評価のしかたは\figref{fig:ml-sgd}と同じ。点は種 3 つの中央値、縦の線は最小から最大まで（多くは点に隠れる）。線の色と印は左右で共通で（図の下の凡例）、青の丸が $B=1$、赤の四角が $B=10$、緑の三角が $B=100$。
左：$(0,0)$ から 3 エポック。右：谷底から 1 エポック。右では、どの線も右上がりで、上から $B=1$、$10$、$100$ の順に並ぶ。}
\label{fig:ml-sgd-eta}
\end{figure}
% [/figpass]
学習の途中で学習率を下げていく工夫は、\chapref{ch:training}で扱います。

バッチサイズは、小さければよいわけでもありません。
計算機は、$B$ 件をまとめて行列の計算にすると、1 件ずつよりずっと速く処理できます（\chapref{ch:pytorch}）。
そのため、メモリに収まる範囲で大きめの $B$ を選びます。
\tabref{tab:ml-scale}は、この章の実験と実際の LLM の比較です。

\begin{table}[tbp]
\centering
\caption{ミニバッチ学習の規模。言語モデルでは、文章の中の 1 か所がデータ 1 件にあたるので、データもバッチサイズもトークン数で数える。
本書の基盤モデルの 4,096 トークンは、128 トークンの列 32 本分（トークンの種類の数 4096 と同じなのは偶然）。
本書の基盤モデルの値は\chapref{ch:pretraining}の設定、LLaMA は\cite{touvron2023llama}、Llama 3.1 は\cite{dubey2024llama3}による。
モデル名の 65B と 405B は、パラメータの数（650 億個と 4,050 億個）を表し、バッチサイズ $B$ とは関係がない。}
\label{tab:ml-scale}
\small
\begin{tabular}{lrrr}
\toprule
 & 訓練データ & バッチサイズ & ステップ数 \\
\midrule
この節の実験 & 1 万件 & 10 件 & 3,000 \\
本書の基盤モデル & 約 250 万トークン & 4,096 トークン & 4,000 \\
LLaMA 65B（2023 年） & 約 1.4 兆トークン & 約 400 万トークン & 約 35 万 \\
Llama 3.1 405B（2024 年） & 約 15.6 兆トークン & 400 万〜1,600 万トークン & 約 120 万 \\
\bottomrule
\end{tabular}
\end{table}

規模は何桁も違いますが、していることは同じです。
なお、本書の基盤モデルはデータを約 6.6 周（エポック）します。
実際の LLM は、データが膨大なので、ほぼ 1 周しかしません。
LLaMA のステップ数は、データ量をバッチサイズで割って求めた概算です。

% [figpass] 案 11
最後に、全データを使う勾配降下法と、$B$ の両極端を含むミニバッチ SGD を、この節の実験の値で\tabref{tab:ml-gd-vs-sgd}に並べます。
$B=N$ と $B=1$ は両端で、ふつうはその間の $B$ を使います。

\begin{table}[tbp]
\centering
\caption{全データの勾配降下法と、ミニバッチ SGD の比較（データ 1 万件、学習率 0.1 の、この節の実験の値）。
勾配のばらつきは $w=b=0$ の地点での $\pd{\calL}{w}$ の標準偏差、損失の最小値は $0.0091$。}
\label{tab:ml-gd-vs-sgd}
\small
\begin{tabularx}{\linewidth}{l>{\raggedright\arraybackslash}X>{\raggedright\arraybackslash}X>{\raggedright\arraybackslash}X}
\toprule
 & 勾配降下法（$B=N$） & SGD（$B=1$） & ミニバッチ SGD（$1<B<N$） \\
\midrule
更新に使うデータ & 全 1 万件 & 1 件 & $B$ 件 \\
1 エポックの更新回数 & 1 & 1 万 & $N/B$（$B=10$ なら 1000） \\
勾配のばらつき & 0（毎回同じ） & 1.446 & $B=10$ で 0.459、$B=100$ で 0.139 \\
1 万件処理後の損失 & 0.167 & 0.0106 & $B=10$ でも $100$ でも 0.0093 \\
谷底のまわりの揺れ & なし & 最も大きい & $B$ が大きいほど小さい（\figref{fig:ml-sgd-eta}） \\
まとめて計算 & できるが、1 回の更新に全データ分の時間とメモリが要る & できない（1 件ずつ） & できる（メモリに収まる $B$ を選ぶ） \\
本書での使いどころ & 小さなデータ（\secref{sec:ml-gd}） & 比較のためだけ & 大きなデータの学習（\tabref{tab:ml-scale}） \\
\bottomrule
\end{tabularx}
\end{table}
% [/figpass]

この節では、確率的勾配降下法を学びました。
全データの平均の勾配を、ステップごとに取り替える一部のデータ（ミニバッチ）の平均の勾配で代用する方法です。
見積もりはぶれますが、同じ計算量で何倍も多く進めるので、大きなデータでは、ほぼ必ず使います。

\section{汎化：見たことのないデータで当てる}\label{sec:ml-generalization}

ここまでは、訓練データでの損失を小さくすることだけを考えてきました。
しかし、本当に当てたいのは、訓練データの正解ではありません。それはすでに手元にあります。
当てたいのは、\textbf{まだ見ていないデータ}の正解です。
この節の問いは、「訓練データで損失が小さければ、良いモデルと言えるのか」です。
4 要素で言えば、データの使い方の話です。

極端なモデルを考えます。
訓練データの 100 件をすべて表に書き写し、入力が表にあれば、その正解を答えるモデルです。
表にない入力には、いつも 0 と答えます。
このモデルの訓練データでの損失は 0 ですが、新しい断片のトークン数は、まったく当てられません。
試験問題と答えを丸暗記した学生が、初めて見る問題を解けないのと同じです。

訓練データに限らず、新しいデータでも正しく予測できることを\term{汎化}{generalization}と呼びます。
汎化の程度を測るには、学習に使わなかったデータで損失を計算します。
そのために、手元のデータを最初に三つに分けておきます。
\term{訓練データ}{training data}は、勾配を計算して、パラメータを更新するために使います。
\term{検証データ}{validation data}は、学習の途中で汎化の程度を測るために使います。
勾配は計算しないので、パラメータの更新には使いません。
その代わり、測った結果を見て、人が設定（多項式の次数や学習率など）を選び、学習のやめどきを決めます。
\term{テストデータ}{test data}は、最後に一度だけ使って、完成したモデルの成績を測ります。
その結果は、どの判断にも使いません。
三つの役割を\figref{fig:ml-data-split}に並べました（検証とテストを分ける理由は、節の最後で述べます）。

\begin{figure}[tbp]
\centering
\begin{tikzpicture}[font=\footnotesize,
  part/.style={minimum height=8mm,inner sep=0pt,rounded corners=0pt},
  use/.style={font=\footnotesize,align=right,inner sep=1.5pt}]
  \node[box,part,minimum width=64mm] (tr) {\textbf{訓練データ}};
  \node[orangebox,part,minimum width=30mm,anchor=west] (va) at ($(tr.east)+(1mm,0)$) {\textbf{検証データ}};
  \node[redbox,part,minimum width=30mm,anchor=west] (te) at ($(va.east)+(1mm,0)$) {\textbf{テストデータ}};
  \node[lbl,anchor=south west] at ($(tr.north west)+(0,0.8mm)$) {手元のデータ（学習の前に、三つに分けておく）};
  % モデルへの影響（訓練は直接、検証は間接、テストは影響しない）
  \node[greenbox,part,minimum width=95mm,anchor=north west] (mdl) at ($(tr.south west)+(0,-15mm)$)
    {\textbf{モデル}（パラメータ $\vtheta$）と、\textbf{学習の設定}（ハイパーパラメータ）};
  \draw[arr] (tr.south) -- node[left,use]{毎ステップ使う。\\勾配でパラメータを\\更新する（直接）} (tr.south |- mdl.north);
  \draw[darr] (va.south) -- node[left,use]{損失を測るだけ（勾配なし）。\\結果を見て、人が設定を選び、\\やめどきを決める（間接）} (va.south |- mdl.north);
  \node[use,align=center,anchor=north] at ($(te.south)+(0,-1.5mm)$) {最後に一度、成績を\\測るだけ。結果は、\\どの判断にも使わない};
\end{tikzpicture}
\caption{データの三つの役割と、モデルへの影響。訓練データは、勾配でパラメータを直接更新する（実線）。
検証データは、パラメータを更新しないが、設定選びとやめどきの判断を通して、できあがるモデルに間接的に影響する（破線）。
テストデータは、成績を測るだけで、モデルにも学習の進め方にも影響しない（モデルへ向かう矢印がない）。
帯の幅は割合の一例。}
\label{fig:ml-data-split}
\end{figure}

\secref{sec:ml-gd}の直線を、学習に使っていない 100 件で評価すると、損失は $0.0076$ でした。
訓練データでの損失 $0.0112$ と同じ程度です\footnote{訓練データのほうが大きいのは、114 字で 120 トークンという極端な断片が 1 件含まれていたためです。
珍しい文字は、1 文字が複数のトークンに分かれることがあります（\chapref{ch:tokenizer}）。}。
この 100 件や、前節までの「学習に使っていない」データは、設定選びにも使っていないので、役割はテストデータにあたります。
パラメータが 2 個しかないモデルは、100 件を丸暗記できないので、新しいデータでも同じように当たります。
では、モデルをもっと柔軟にすると、何が起きるでしょうか。

\subsection*{多項式回帰で過学習を見る}

正解の曲線が分かっている人工データで実験します。
$y=\sin(\pi x)$ に、小さな偶然のずれを足した点を、10 件だけ訓練データにします。
モデルは $M$ 次の多項式です。
\begin{equation}
f_\theta(x) = w_0 + w_1 x + w_2 x^2 + \dots + w_M x^M
\label{eq:ml-poly}
\end{equation}
$M$ を\textbf{次数}と呼びます。$M=1$ なら直線で、$M$ が大きいほど、曲がりくねった曲線を表せます。
パラメータは $w_0,\dots,w_M$ の $M+1$ 個です。
$x^2$ や $x^3$ を「別の入力」と見れば、これは入力が複数ある線形回帰（\eqref{eq:ml-linear-multi}）です。
グラフは曲線ですが、予測がパラメータ（重み）の 1 次式なので、「線形」と呼びます。
損失は二乗誤差です。
線形回帰なので、\secref{sec:ml-gd}の終わりで触れた公式で、損失が最小になる係数を直接求めました（\file{code/ch03/overfit_poly.py} の \code{fit} 関数）。
勾配降下法の止めどきに左右されずに、次数の違いだけを比べるためです。
次数を 0 から 9 まで変えて当てはめました。奇数の次数の結果を抜き出したものが次です。
誤差は、$y$ と同じ単位で読めるように、平均二乗誤差の平方根で表します。
検証誤差は、同じ曲線から作った別の 200 件で測りました。

\outputfile{ch03/out/poly_degree_odd.txt}

訓練誤差は、次数を上げるほど下がり、$M=9$ で 0 になります。
パラメータが 10 個あれば、10 個の点（$x$ の値はすべて違います）をすべて通る曲線を作れるからです。
点を 1 つ通るごとに条件が 1 つ増えますが、調整できる数も同じ 10 個あるので、すべての条件を満たせます。
ところが、検証誤差は $M=3$ の $0.290$ が最小です。
そのあとは $M=7$ までほとんど変わらず、$M=8$ と $9$ で大きく上がります（全部の次数は\figref{fig:ml-poly-errors}の左）。
\figref{fig:ml-poly-fits}で、3 つの次数の曲線を見比べてください。

\begin{figure}[tbp]
\centering
\pgfplotsset{mlpoly/.style={width=0.385\linewidth,height=0.29\linewidth,
  xmin=-1.05,xmax=1.05,ymin=-1.9,ymax=1.9,xtick={-1,0,1},ytick={-1,0,1},xlabel={$x$},
  tick label style={font=\scriptsize},label style={font=\footnotesize},title style={font=\footnotesize,yshift=-1.5mm}}}
\begin{tikzpicture}
\begin{axis}[mlpoly,title={$M=1$：未学習},ylabel={$y$},ylabel style={yshift=-2mm}]
\addplot[bookgreen,dashed,thick] table[col sep=comma,x=x,y=true]{code/ch03/out/poly_curves.csv};
\addplot[bookaccent,very thick] table[col sep=comma,x=x,y=m1]{code/ch03/out/poly_curves.csv};
\addplot[only marks,mark=*,mark size=1.5pt,bookblue] table[col sep=comma,x=x,y=y]{code/ch03/out/poly_train.csv};
\end{axis}
\end{tikzpicture}\hfill
\begin{tikzpicture}
\begin{axis}[mlpoly,title={$M=3$：ちょうどよい},yticklabels={}]
\addplot[bookgreen,dashed,thick] table[col sep=comma,x=x,y=true]{code/ch03/out/poly_curves.csv};
\addplot[bookaccent,very thick] table[col sep=comma,x=x,y=m3]{code/ch03/out/poly_curves.csv};
\addplot[only marks,mark=*,mark size=1.5pt,bookblue] table[col sep=comma,x=x,y=y]{code/ch03/out/poly_train.csv};
\end{axis}
\end{tikzpicture}\hfill
\begin{tikzpicture}
\begin{axis}[mlpoly,title={$M=9$：過学習},yticklabels={}]
\addplot[bookgreen,dashed,thick] table[col sep=comma,x=x,y=true]{code/ch03/out/poly_curves.csv};
\addplot[bookaccent,very thick] table[col sep=comma,x=x,y=m9]{code/ch03/out/poly_curves.csv};
\addplot[only marks,mark=*,mark size=1.5pt,bookblue] table[col sep=comma,x=x,y=y]{code/ch03/out/poly_train.csv};
\end{axis}
\end{tikzpicture}
\caption{10 件の訓練データ（青い点）に、次数の違う多項式（赤い線）を当てはめた。緑の破線が、データを作った本当の曲線。
直線は単純すぎて、曲線を表せない。9 次式は 10 個の点をすべて通るが、点と点の間で上下に大きく振れる（とくに左端の近く）。}
\label{fig:ml-poly-fits}
\end{figure}

$M=1$ の直線は、訓練データにも合っていません。
このように、訓練データにさえ合っていない状態を\term{未学習}{underfitting}と呼びます（「まだ学習していない」という意味ではありません）。
ここでの原因は、モデルが単純すぎて、データの形を表せないことです。
$M=3$ は、本当の曲線（緑の破線）に近い形です。
$M=9$ は、10 個の点を正確に通ります。
しかし、点と点の間や両端で、大きく波打っています。
訓練データに含まれる偶然のずれまで、忠実になぞってしまったのです。
訓練データには合うのに、新しいデータで外れる状態を\term{過学習}{overfitting}と呼びます。
丸暗記のモデルは、過学習の極端な例でした。

\figref{fig:ml-poly-errors}の左は、次数と誤差の関係です。
訓練誤差（青）は下がり続けます。
検証誤差（赤）は、$M=3$ で最小になったあとは下がらなくなり、$M=8$ から大きく上がります。
訓練誤差は下がり続けるのに、検証誤差が下がらなくなり、やがて上がり始めたら、過学習のしるしです。
二つの誤差に差があること自体は、ふつうのことです（ちょうどよい $M=3$ でも、訓練誤差 $0.215$ に対して、検証誤差は $0.290$ です）。
学習中に訓練損失と検証損失の両方を見るのは、このためです。
本書の基盤モデルでは、学習の終わりの訓練損失が $3.11$、検証損失が $3.48$ でした。
この差（$0.37$）だけでは過学習とは言えず、検証損失が下がり続けているかどうかを合わせて見ます（\chapref{ch:pretraining}で実際に確かめます）。

\begin{figure}[tbp]
\centering
\begin{tikzpicture}
\begin{axis}[width=0.5\linewidth,height=0.315\linewidth,xlabel={次数 $M$},ylabel={誤差},
  xmin=-0.3,xmax=9.3,ymin=0,ymax=0.85,xtick={0,1,...,9},
  tick label style={font=\scriptsize},label style={font=\footnotesize},
  legend columns=2,
  legend style={font=\scriptsize,at={(0.98,0.97)},anchor=north east,draw=black!30,/tikz/every even column/.append style={column sep=1.5mm}},legend cell align=left]
\addplot[bookblue,thick,mark=*,mark size=1.4pt] table[col sep=comma,x=M,y=train]{code/ch03/out/poly_errors.csv};
\addplot[bookaccent,thick,mark=square*,mark size=1.4pt] table[col sep=comma,x=M,y=val]{code/ch03/out/poly_errors.csv};
\legend{訓練誤差,検証誤差}
\node[font=\scriptsize,anchor=north] at (axis cs:1.4,0.52) {未学習};
\node[font=\scriptsize,anchor=south] at (axis cs:8.2,0.47) {過学習};
\end{axis}
\end{tikzpicture}\hfill
\begin{tikzpicture}
\begin{axis}[width=0.5\linewidth,height=0.315\linewidth,xlabel={$x$},ylabel={$y$},
  xmin=-1.05,xmax=1.05,ymin=-1.9,ymax=1.9,
  tick label style={font=\scriptsize},label style={font=\footnotesize},
  legend style={font=\scriptsize,at={(0.03,0.97)},anchor=north west,draw=black!30},legend cell align=left]
\addplot[black!45,thin] table[col sep=comma,x=x,y=m9]{code/ch03/out/poly_curves.csv};
\addplot[bookaccent,very thick] table[col sep=comma,x=x,y=reg_good]{code/ch03/out/poly_curves.csv};
\addplot[bookgreen,dashed,thick] table[col sep=comma,x=x,y=true]{code/ch03/out/poly_curves.csv};
\addplot[only marks,mark=*,mark size=1.5pt,bookblue] table[col sep=comma,x=x,y=y]{code/ch03/out/poly_train.csv};
\legend{$\lambda=0$,$\lambda=0.001$}
\end{axis}
\end{tikzpicture}
\caption{左：多項式の次数と誤差。訓練誤差は下がり続ける。検証誤差は $M=3$ で最小になり、$M=7$ まではほぼ横ばいで、$M=8$ から急に上がる。
右：同じ 9 次式に L2 正則化をかけた結果（赤）。正則化なし（灰色）の大きな振れが収まり、本当の曲線（緑の破線）に近づく。青い点は訓練データ。}
\label{fig:ml-poly-errors}
\end{figure}


\subsection*{過学習を防ぐ}
\skimread{「対策は三つある」という結論だけ押さえれば、先へ進めます。データを増やす、モデルを単純にする、重みが大きくなりすぎないように、重みの 2 乗の和を損失に足す（L2 正則化）、の三つです。
足す項の強さを決める数を $\lambda$ と書きます。}

過学習への対策は、大きく三つあります。

\textbf{第一に、データを増やします。}
9 次式のまま、訓練データを 10 件から 100 件に増やすと、検証誤差は $0.459$ から $0.283$ に下がりました。
1000 件では $0.275$ です。
データが多ければ、偶然のずれは打ち消し合い、柔軟なモデルでも本当の形を学べます。
LLM が膨大な文章で学習するのは、巨大なモデルを過学習させないためでもあります。
% [figpass] 案 9
\figref{fig:ml-poly-sweeps}の左下は、\file{code/ch03/sweep_poly.py} で訓練データの件数を 10〜1000 件の 6 通りに変え、データを作る乱数の種も 10 通り変えて、3 次式と 9 次式の検証誤差を比べたものです。
10 件では、9 次式の検証誤差は種によって 0.42 から 1.45 までばらつき、3 次式よりずっと悪くなります。
件数が増えるほど 2 本の差は縮まり、300 件以上では、9 次式も 3 次式と同じ程度になります。
どちらも、本当の曲線そのものの検証誤差（0.270。偶然のずれの分）より下には行きません。

\begin{figure}[tbp]
\centering
\pgfplotsset{mllam/.style={width=0.3\linewidth,height=0.25\linewidth,
  xmin=-1.05,xmax=1.05,ymin=-1.9,ymax=1.9,xtick={-1,0,1},ytick={-1,0,1},
  tick label style={font=\scriptsize},label style={font=\scriptsize},title style={font=\footnotesize,yshift=-1.5mm}}}
\begin{tikzpicture}
\begin{axis}[mllam,title={$\lambda=0$},ylabel={$y$},ylabel style={yshift=-2mm}]
\addplot[bookgreen,dashed,thick] table[col sep=comma,x=x,y=true]{code/ch03/out/sweep_poly_curves.csv};
\addplot[bookaccent,very thick] table[col sep=comma,x=x,y=lam0]{code/ch03/out/sweep_poly_curves.csv};
\addplot[only marks,mark=*,mark size=1.2pt,bookblue] table[col sep=comma,x=x,y=y]{code/ch03/out/poly_train.csv};
\end{axis}
\end{tikzpicture}\hfill
\begin{tikzpicture}
\begin{axis}[mllam,title={$\lambda=10^{-5}$},yticklabels={}]
\addplot[bookgreen,dashed,thick] table[col sep=comma,x=x,y=true]{code/ch03/out/sweep_poly_curves.csv};
\addplot[bookaccent,very thick] table[col sep=comma,x=x,y=lam1e-5]{code/ch03/out/sweep_poly_curves.csv};
\addplot[only marks,mark=*,mark size=1.2pt,bookblue] table[col sep=comma,x=x,y=y]{code/ch03/out/poly_train.csv};
\end{axis}
\end{tikzpicture}\hfill
\begin{tikzpicture}
\begin{axis}[mllam,title={$\lambda=10^{-3}$},yticklabels={}]
\addplot[bookgreen,dashed,thick] table[col sep=comma,x=x,y=true]{code/ch03/out/sweep_poly_curves.csv};
\addplot[bookaccent,very thick] table[col sep=comma,x=x,y=lam1e-3]{code/ch03/out/sweep_poly_curves.csv};
\addplot[only marks,mark=*,mark size=1.2pt,bookblue] table[col sep=comma,x=x,y=y]{code/ch03/out/poly_train.csv};
\end{axis}
\end{tikzpicture}\hfill
\begin{tikzpicture}
\begin{axis}[mllam,title={$\lambda=1$},yticklabels={}]
\addplot[bookgreen,dashed,thick] table[col sep=comma,x=x,y=true]{code/ch03/out/sweep_poly_curves.csv};
\addplot[bookaccent,very thick] table[col sep=comma,x=x,y=lam1]{code/ch03/out/sweep_poly_curves.csv};
\addplot[only marks,mark=*,mark size=1.2pt,bookblue] table[col sep=comma,x=x,y=y]{code/ch03/out/poly_train.csv};
\end{axis}
\end{tikzpicture}

\vspace{1mm}
\begin{tikzpicture}
\begin{semilogxaxis}[width=0.52\linewidth,height=0.33\linewidth,xlabel={訓練データの件数},ylabel={検証誤差},
  xmin=7,xmax=1500,ymin=0.2,ymax=1.5,
  tick label style={font=\scriptsize},label style={font=\footnotesize},
  legend style={font=\scriptsize,at={(0.97,0.97)},anchor=north east,draw=black!30},legend cell align=left]
\addplot[bookaccent,thick,mark=square*,mark size=1.3pt,error bars/.cd,y dir=both,y explicit]
  table[col sep=comma,x=n,y=median,y error plus expr=\thisrow{hi}-\thisrow{median},y error minus expr=\thisrow{median}-\thisrow{lo}]{code/ch03/out/sweep_poly_n_m9.csv};
\addplot[bookblue,thick,mark=*,mark size=1.3pt,error bars/.cd,y dir=both,y explicit]
  table[col sep=comma,x=n,y=median,y error plus expr=\thisrow{hi}-\thisrow{median},y error minus expr=\thisrow{median}-\thisrow{lo}]{code/ch03/out/sweep_poly_n_m3.csv};
\addplot[bookgreen,thick,dashed] table[col sep=comma,x=x,y=true]{code/ch03/out/sweep_poly_lam0.csv};
\legend{9 次式,3 次式,本当の曲線}
\end{semilogxaxis}
\end{tikzpicture}\hfill
\begin{tikzpicture}
\begin{semilogxaxis}[width=0.5\linewidth,height=0.33\linewidth,xlabel={正則化の強さ $\lambda$},ylabel={誤差},
  xmin=3e-9,xmax=30,ymin=-0.03,ymax=1.2,xtick={1e-8,1e-6,1e-4,1e-2,1},ytick={0,0.2,0.4,0.6,0.8},
  tick label style={font=\scriptsize},label style={font=\footnotesize},
  legend style={font=\scriptsize,at={(0.03,0.97)},anchor=north west,draw=black!30},legend cell align=left]
\addplot[bookblue,thick,mark=*,mark size=1pt] table[col sep=comma,x=lam,y=train]{code/ch03/out/sweep_poly_lam.csv};
\addplot[bookaccent,thick,mark=square*,mark size=1pt] table[col sep=comma,x=lam,y=val]{code/ch03/out/sweep_poly_lam.csv};
\addplot[black!45,thin,densely dashed] table[col sep=comma,x=x,y=val0]{code/ch03/out/sweep_poly_lam0.csv};
\legend{訓練誤差,検証誤差,$\lambda=0$ の検証誤差}
\end{semilogxaxis}
\end{tikzpicture}
\caption{過学習の対策を、パラメータを変えて見る（\file{code/ch03/sweep_poly.py}）。
上：9 次式に、強さ $\lambda$ を変えた L2 正則化をかけた曲線（赤）。青い点は訓練データ 10 件、緑の破線は本当の曲線。
左下：訓練データの件数と検証誤差。点はデータの乱数の種 10 通りの中央値、縦の線は最小から最大まで。
緑の破線は、本当の曲線そのものの検証誤差。右下：9 次式（訓練データ 10 件）で $\lambda$ を 19 通りに変えた誤差（横軸は対数目盛り）。
灰色の破線は正則化なし（$\lambda=0$）の検証誤差 0.459。検証データは、どれも本文と同じ 200 件。}
\label{fig:ml-poly-sweeps}
\end{figure}
% [/figpass]

\textbf{第二に、モデルを単純にします。}
次数を 9 から 3 に下げるのが、その例です。

\textbf{第三に、パラメータが大きくなりすぎないようにします。}
先ほどの出力の右端の列を見てください。
$M=9$ では、係数の絶対値が最大で $136.8$ にもなっています。
曲線が大きく振れるのは、絶対値の大きな正の項と負の項を足し合わせて、点を通しているからです。
点から少し離れると、この打ち消し合いがずれて、値が大きく振れます。
% [figpass] 案 10
\figref{fig:ml-poly-terms}は、その様子を数値で見たものです（\file{code/ch03/sweep_poly.py} の出力）。
訓練データの点 $x=-0.778$ では、絶対値が 20 を超える項どうしが打ち消し合い、和は正解の $-0.396$ に一致します。
隣の点との中間 $x=-0.889$ では、定数項を除く各項が 1.1〜3.3 倍に（次数が高い項ほど大きく）なり、打ち消し合いがずれて、和は $-1.603$ です。
本当の曲線の値は $-0.342$ なので、1 以上も外れています。

\begin{figure}[tbp]
\centering
\begin{tikzpicture}
\begin{axis}[width=0.8\linewidth,height=0.32\linewidth,ybar,bar width=4.2pt,
  xlabel={次数 $j$（項 $w_jx^j$ の番号）},ylabel={項の値 $w_jx^j$},
  xmin=-0.6,xmax=9.6,xtick={0,1,...,9},ymin=-65,ymax=52,
  tick label style={font=\scriptsize},label style={font=\footnotesize},
  legend style={font=\scriptsize,at={(0.5,1.03)},anchor=south,draw=black!30,legend columns=2,/tikz/every even column/.append style={column sep=2mm}},legend cell align=left]
\addplot[fill=figblue,draw=bookblue] table[col sep=comma,x=j,y=on]{code/ch03/out/sweep_poly_terms.csv};
\addplot[fill=figorange,draw=bookorange] table[col sep=comma,x=j,y=mid]{code/ch03/out/sweep_poly_terms.csv};
\legend{$x=-0.778$（訓練データの点）：和 $-0.396$,$x=-0.889$（中間）：和 $-1.603$}
\draw[black!50,thin] (axis cs:-0.6,0) -- (axis cs:9.6,0);
\end{axis}
\end{tikzpicture}
\caption{正則化なしの 9 次式（\figref{fig:ml-poly-fits}の右）の 10 個の項 $w_jx^j$ を、2 つの $x$ で並べた棒グラフ。
予測は 10 本の棒の和。訓練データの点（青）では、大きな正と負の項が打ち消し合って正解に一致する。
中間の点（橙）では、どの項もさらに大きくなり、和は本当の曲線（$-0.342$）から 1 以上ずれる。}
\label{fig:ml-poly-terms}
\end{figure}
% [/figpass]
そこで、損失に、パラメータの大きさを表す項を足します。
この項を\textbf{ペナルティ}（罰則）の項、または\textbf{正則化項}と呼びます。
\begin{equation}
\calL_{\text{正則化}}(\vtheta) = \calL(\vtheta) + \lambda \norm{\vw}^2,
\qquad \norm{\vw}^2 = w_1^2 + w_2^2 + \dots + w_M^2
\label{eq:ml-l2}
\end{equation}
第 2 項は、重みの 2 乗の和です。
定数項 $w_0$ は曲線全体を上下にずらすだけなので、ペナルティには入れません。$\vw$ は、$w_0$ 以外の重みを並べたベクトルです。
重みが大きいほど損失が増えるので、最適化は「データに合わせつつ、重みは小さく」を目指します。
この工夫を \textbf{L2 正則化}（L2 regularization）と呼びます。
$\lambda$（ラムダ）は、ペナルティの強さを決める数です。

9 次式に、$\lambda$ を変えながら L2 正則化をかけてみました。
$\lambda=0.001$ では、係数の最大値は $136.8$ から $2.8$ に収まり、検証誤差は $0.459$ から $0.286$ に下がりました。
次数 3 の多項式と同じ程度の成績です（\figref{fig:ml-poly-errors}の右）。
一方、$\lambda=1$ ではペナルティが強すぎて、訓練誤差が $0.645$、検証誤差が $0.700$ となり、未学習になりました。
深層学習では、L2 正則化とほぼ同じ働きをする\textbf{重み減衰}（weight decay）がよく使われます（\chapref{ch:training}）。
二つの違いも、そこで説明します。

% [figpass] 案 8
\figref{fig:ml-poly-sweeps}の上の 4 枚は、$\lambda$ を 4 通りに変えた 9 次式の曲線を並べたものです。
$\lambda$ を大きくするほど振れが収まり、$\lambda=1$ では曲線がほぼ平らになります（未学習）。
右下は、$\lambda$ を $10^{-8}$ から $10$ まで、10 の 0.5 乗倍ずつ 19 通りに変えた誤差です。
訓練誤差は $\lambda$ とともに上がり続けますが、検証誤差は $10^{-5}$〜$10^{-3}$ のあたりで低くなり、その両側で上がります。
$\lambda$ が小さすぎれば過学習、大きすぎれば未学習で、次数 $M$ を変えたとき（\figref{fig:ml-poly-errors}の左）と、左右を入れ替えた形です。


三つの対策を\tabref{tab:ml-overfit-remedies}にまとめます。
どの対策にも「行き過ぎ」があり（データを増やす場合は、その手間）、どこで止めるかは検証データで決めます。

\begin{table}[tbp]
\centering
\caption{過学習の三つの対策。この章の結果は、9 次式・訓練データ 10 件（検証誤差 0.459）から始めたときの検証誤差。}
\label{tab:ml-overfit-remedies}
\small
\begin{tabularx}{\linewidth}{l>{\raggedright\arraybackslash}Xl>{\raggedright\arraybackslash}X>{\raggedright\arraybackslash}X}
\toprule
対策 & 変えるもの & 選ぶ値 & この章の結果 & 行き過ぎると・費用 \\
\midrule
データを増やす & 訓練データの件数 & --- & 100 件で 0.283、1000 件で 0.275 & 集める手間と、学習の計算が増える \\
モデルを単純にする & モデルの形 & 次数 $M$ & $M=3$ で 0.290 & 単純すぎると未学習（$M=1$ で 0.564） \\
L2 正則化 & 損失（ペナルティの項を足す） & 強さ $\lambda$ & $\lambda=0.001$ で 0.286 & 強すぎると未学習（$\lambda=1$ で 0.700） \\
\bottomrule
\end{tabularx}
\end{table}
% [/figpass]

\subsection*{ハイパーパラメータと、データを三つに分ける理由}

多項式の次数 $M$ や、正則化の強さ $\lambda$ は、勾配降下法では決まりません。
学習の前に、人が決める値です。
このような値を\term{ハイパーパラメータ}{hyperparameter}と呼びます（ここまで「設定」と呼んできたものです）。
学習率 $\eta$、バッチサイズ $B$、ステップ数もハイパーパラメータです。
パラメータは訓練データで学習し、ハイパーパラメータは検証データで選びます。
モデルの柔軟さを決める設定（$M$ や $\lambda$）を訓練誤差で選ぶと、いつも最も柔軟なもの（$M=9$、正則化なしの $\lambda=0$）が選ばれてしまうからです。
学習率も、最後は検証データで選びます（\secref{sec:ml-gd}で訓練データでの損失の下がり方を見たのは、発散しない範囲に絞るためです）。

ここで、検証データとテストデータの違いを整理します（\figref{fig:ml-data-split}）。
どちらも、勾配を計算しないので、パラメータの更新には使いません。
違いは、測った結果を、何かを決めることに使うかどうかです。
検証データの結果は、ハイパーパラメータ選びと、学習のやめどきの判断に使います。
つまり、検証データは、パラメータを直接は更新しませんが、これらの判断を通して、できあがるモデルに間接的に影響します。
テストデータの結果は、どの判断にも使わないので、モデルにも学習の進め方にも影響しません。

この違いから、テストデータが必要な理由が分かります。
いくつもの設定の中から検証誤差が最小のものを選ぶと、その検証データにたまたま合っていた設定を選びやすくなります。
そのため、選ばれた設定の検証誤差は、モデルの実力より少し良く見える値になりがちです。
まだ見ていないデータでの実力を見積もるには、選ぶことに一度も使っていないテストデータが必要なのです。
先ほどの実験では、検証誤差が最小の $M=3$ が選ばれ、その検証誤差は $0.290$ でした。
この多項式を、さらに別の 200 件のテストデータで測ると、誤差は $0.332$ でした。
差には偶然も含まれますが、成績として報告するのは、選ぶことに使った $0.290$ ではなく、この $0.332$ です。
テストデータを見ながら設定を直してはいけません。
テストデータも設定選びに使ったことになり、検証データと同じ役割になってしまうからです（入試の本番の問題で練習するのと同じです）。

この節では、訓練データでの損失と新しいデータでの損失は別物だと学びました。
柔軟すぎれば過学習、単純すぎれば未学習です。
釣り合いは検証データで確かめ、最後の成績は、どの判断にも使っていないテストデータで測ります。

\section{学習の種類を整理する}\label{sec:ml-map}
\skimread{\tabref{tab:ml-map}を眺めて、「言語モデルの事前学習は、文章自身から正解を作る自己教師あり学習である」という 1 点だけ押さえれば十分です。
\chapref{ch:lm}で、もう一度説明します。}

最後に、この章で学んだ枠組みを、本書全体の中に位置づけます。
この章の問題には、どれも、入力とは別に用意した正解がありました（トークン数はトークナイザで数え、足し算の答えは計算ドリルから取りました）。
これが教師あり学習でした（正解は、人が 1 件ずつ付けることもあります）。
機械学習には、正解の与え方が違う種類の学習もあります（\tabref{tab:ml-map}）。

\begin{table}[tbp]
\centering
\caption{学習の種類。違いは「何を手がかりに学ぶか」にある。4 要素の枠組みは共通である。}
\label{tab:ml-map}
\small
\begin{tabularx}{\linewidth}{llX}
\toprule
種類 & 学習の手がかり & 例と、本書で扱う場所 \\
\midrule
教師あり学習 & 入力とは別に用意した正解 & この章、事後学習（\chapref{ch:posttraining}） \\
教師なし学習 & なし（データどうしの類似だけ） & 似たデータをまとめる（本書の範囲外） \\
自己教師あり学習 & 入力データ自身の一部 & 事前学習（\chapref{ch:lm}、\chapref{ch:pretraining}） \\
強化学習 & 良し悪しを表す数（報酬） & 迷路でゴールに着く学習（\chapref{ch:rl-basics}）、正解なら報酬 1 として計算問題の正答率を上げる学習（\chapref{ch:rl-llm}） \\
\bottomrule
\end{tabularx}
\end{table}

\textbf{強化学習}（reinforcement learning）では、正解そのものは与えられず、答えの良し悪しを表す数（報酬）だけが返ってきます。
LLM にとって最も重要なのは、\textbf{自己教師あり学習}（self-supervised learning）です。
人が正解を付ける方法では、兆の単位のデータは用意できません。
しかし、文章があれば、「途中まで」を入力、「その次のトークン」を正解とする組を、いくらでも自動で作れます。
正解を、入力と同じデータ（文章）の中から取り出すところが、教師あり学習との違いです。
この章の softmax 回帰と違うのは、正解の作り方と、ロジットを作る部分です（次章から、$\mW\vx+\vb$ より表現力の高いものに取り替えます）。
softmax で確率に直すこと、交差エントロピー損失、勾配を使った最適化は、この章と同じ形です。

\begin{summary}
\item 機械学習は、規則を書く代わりに、パラメータを持つ関数（モデル）を、例に合うように調整します。
枠組みは、データ・モデル・損失関数・最適化の 4 要素です。
教師あり学習では、訓練データでの損失の平均を最小にします（経験リスク最小化）。本当に小さくしたいのは、まだ見ていないデータでの損失です。
\item 勾配降下法は、勾配の逆向きにパラメータを少しずつ動かします。学習率が小さすぎると遅く、大きすぎると発散します。
\item 損失は、正解についての確率分布を出力するようにモデルを作れば、その負の対数尤度として作れます。分布に正規分布を選ぶと、二乗誤差になります。
\item 分類では、ロジットを softmax で確率に直し、交差エントロピー損失を使います。ロジットについての勾配は $\vp-\vy$ です。
\item 確率的勾配降下法は、全データの平均の勾配を、毎回取り替えるミニバッチで見積もります。ぶれますが、偏りはなく、同じ計算量でずっと速く進みます。
\item 訓練データに合うことと、新しいデータで当たること（汎化）は別です。検証データで過学習を見張り、ハイパーパラメータを選びます。テストデータは、最後に成績を測るためだけに使い、どの判断にも使いません。
\end{summary}

\section*{次の章へ}

この章のモデルは、どれも「入力に重みを掛けて足す」形でした。
この形には限界があります。
\figref{fig:ml-blobs}の右で見たとおり、入力をそのまま使う softmax 回帰は直線の境界しか作れず、渦巻きのデータでは正解率が 52.7\% にとどまりました。
言葉の規則は、渦巻きよりはるかに込み入っています。曲がった境界を作れるモデルが必要です。
次の\chapref{ch:nn}では、線形モデルを何層も積み重ねた\textbf{ニューラルネットワーク}を作ります。
\secref{sec:ml-autograd}の 4 手順のループはそのまま使い（次章では、5 行のコードに整理します）、勾配を自動で求める仕組みを、自分の手で実装します。
なお、次章では渦巻きのデータを作り直すので、線形モデルの正解率は 54\% と、少し違う値になります。
半分近くを間違えることは同じです。

\section*{さらに学ぶために}

\begin{itemize}
\item \cite{bishop2006prml}：1 章と 3〜4 章。多項式回帰の過学習と、最尤推定から損失を導く流れの原典です。
\item \cite{goodfellow2016deep}：5 章。汎化、最尤推定、確率的勾配降下法を整理しています。
\item \cite{zhang2023d2l}：線形回帰と softmax 回帰を PyTorch で一から実装する、無料の教科書です。
\item \cite{cs229notes}：スタンフォード大学の講義ノート。線形回帰とロジスティック回帰を数式で丁寧に導きます。
\end{itemize}

\begin{exercises}
\exfirst{1 と 3（どちらも手計算）}
\item \textbf{（手計算）}\ 3 件のデータ $(1,2)$、$(2,3)$、$(3,5)$ に、モデル $f(x)=wx$ を当てはめます。
損失 $\calL(w)$ を $w$ の 2 次式で表し、$w=0$ から学習率 $0.05$ で勾配降下法を 2 ステップ進めてください。
\item\label{ex:ml-scale} \textbf{（コード）}\ \file{code/ch03/linreg_numpy.py} で、\code{x} を作る行の \code{/ 100} を消し、入力 $x$ だけを文字数のままにして学習してください（正解 $y$ は百トークン単位のまま）。
学習率をいくつまで下げれば発散しなくなりますか。
\item \textbf{（手計算）}\ ロジットが $\vz=(1,2,3)$ で、正解がクラス 1（$\vy=(1,0,0)$）のとき、
softmax の出力 $\vp$、交差エントロピー損失、勾配 $\vp-\vy$ を計算してください。
\item\label{ex:ml-sigmoid} \textbf{（導出）}\ (a) $\sigmoid'(z)=\sigmoid(z)\bigl(1-\sigmoid(z)\bigr)$ を示してください。
(b) $K=2$ の softmax で、$p_1 = \sigmoid(z_1 - z_2)$ となることを示してください。
\item\label{ex:ml-laplace} \textbf{（導出）}\ ずれ $\varepsilon$ の確率密度が $\frac{1}{2s}\exp(-|\varepsilon|/s)$（ラプラス分布）だと仮定します。
このとき、負の対数尤度を最小にすることは、どんな損失を最小にすることと同じですか。
\item \textbf{（コード）}\ \file{code/ch03/overfit_poly.py} の訓練データを 20 件に増やすと、検証誤差が最小になる次数は変わりますか。
乱数の種を変えて何度か試し、高い次数の検証誤差がどう変わるかも見てください。
\end{exercises}
